\documentclass[a4paper,11pt,openany]{ltjsbook}
\usepackage{luatexja}
\usepackage{luatexja-fontspec}
\usepackage{amsmath,amssymb,amsfonts,amsthm}
\usepackage{geometry}
\geometry{top=24mm,bottom=24mm,left=24mm,right=24mm,footskip=10mm}
\usepackage{tikz}
\usetikzlibrary{shapes,arrows,positioning,calc,patterns,backgrounds}
\usepackage{tcolorbox}
\tcbuselibrary{skins,breakable}
\usepackage{hyperref}
\usepackage{enumitem}
\usepackage{booktabs}
\usepackage{fancyhdr}

\definecolor{primaryblue}{RGB}{24, 76, 120}
\definecolor{softblue}{RGB}{235, 243, 250}
\definecolor{darkslate}{RGB}{44, 62, 80}
\definecolor{warmamber}{RGB}{217, 119, 6}
\definecolor{softamber}{RGB}{254, 243, 199}
\definecolor{crimson}{RGB}{185, 28, 28}
\definecolor{softrose}{RGB}{255, 228, 230}
\definecolor{accentgreen}{RGB}{16, 120, 70}
\definecolor{softgreen}{RGB}{236, 253, 245}

\hypersetup{
    colorlinks=true,
    linkcolor=primaryblue,
    urlcolor=primaryblue,
    citecolor=accentgreen,
    pdftitle={高校生のための微積分入門 —— 『Calculus Made Easy』で直観的にわかる微分・積分・極限から物理・曲線の数理まで},
    pdfauthor={教育情報分析研究会}
}

\newtcolorbox{examplebox}[1]{
    colback=softgreen,
    colframe=accentgreen,
    coltitle=white,
    fonttitle=\bfseries,
    title=#1,
    arc=3mm,
    boxrule=1pt,
    left=4mm,right=4mm,top=3mm,bottom=3mm,
    breakable
}

\newtcolorbox{pointbox}[1]{
    colback=softamber,
    colframe=warmamber,
    coltitle=darkslate,
    fonttitle=\bfseries,
    title=#1,
    arc=2mm,
    boxrule=1pt,
    left=3mm,right=3mm,top=2mm,bottom=2mm,
    breakable
}

\pagestyle{fancy}
\fancyhf{}
\fancyhead[LE,RO]{\thepage}
\fancyhead[LO]{\nouppercase{\rightmark}}
\fancyhead[RE]{\nouppercase{\leftmark}}
\renewcommand{\headrulewidth}{0.4pt}

\begin{document}

\begin{titlepage}
\begin{center}
\vspace*{1.5cm}
{\Large\bfseries プロジェクト・グーテンベルク古典数学入門シリーズ 第5弾}\\[0.8cm]
\rule{\textwidth}{1.5pt}\\[0.4cm]
{\huge\bfseries 高校生のための微積分入門}\\[0.4cm]
{\Large\bfseries —— 『Calculus Made Easy』で直観的にわかる\\微分・積分・極限から物理・曲線の数理まで ——}\\[0.4cm]
\rule{\textwidth}{1.5pt}\\[1.5cm]

{\large 原著：シルバヌス・P・トンプソン（\textit{Silvanus P. Thompson}）\\
\textit{Calculus Made Easy}（1910年）}\\[2.0cm]

{\Large\bfseries 著：教育情報分析研究会}\\[2.5cm]

\vfill
{\large 2026年10月 最新版}\\[0.5cm]
{\footnotesize 本書は Project Gutenberg のパブリックドメイン数学古典（eBook \#33283）を底本とし、\\
現代の高校数学（数学II・数学III、物理基礎）に完全対応した50ページ以上の本格的入門書です。}
\end{center}
\end{titlepage}

\frontmatter
\tableofcontents

\mainmatter
\chapter{予備的恐怖からの解放：微小量 $dx$ と $dy$ の正体と微積分の扉}
\chaptermark{予備的恐怖からの解放}

\begin{flushright}
\small\textit{To Deliver You From the Preliminary Terrors}
\end{flushright}

\begin{figure}[htbp]
\centering
\begin{tikzpicture}[x=1cm, y=1cm, >=stealth]
    % Outer container
    \draw[thick, fill=softblue, draw=primaryblue, rounded corners=6pt] (-7.6, -4.2) rectangle (7.6, 4.0);
    \node[draw=primaryblue, fill=white, rounded corners=4pt, inner sep=5pt] at (0, 3.2) {\bfseries シルバヌス・P・トンプソン『Calculus Made Easy』の基本精神};
    
    % Differentiation box
    \node[draw=primaryblue, fill=white, rounded corners=4pt, align=center, text width=3.8cm, inner sep=6pt] at (-4.7, 0.9) {
        \textbf{微分（$d$：細分化）}\\[3pt]
        \footnotesize ラテン語 \textit{differentia}\\[2pt]
        \footnotesize ほんの少しの微小増分 $dx$\\[2pt]
        \footnotesize 「微小なかけら」への分解
    };
    
    % Integration box
    \node[draw=crimson, fill=softrose, rounded corners=4pt, align=center, text width=3.8cm, inner sep=6pt] at (4.7, 0.9) {
        \textbf{積分（$\int$：総和化）}\\[3pt]
        \footnotesize ラテン語 \textit{summa}（長いS）\\[2pt]
        \footnotesize 微小量 $f(x)dx$ の連続和\\[2pt]
        \footnotesize 「全体」の組み立て・復元
    };
    
    % Center arrows with clear spacing
    \draw[->, very thick, warmamber] (-2.2, 1.4) -- (2.2, 1.4);
    \node[above=3pt] at (0, 1.4) {\bfseries\small 微分と積分は互いに逆演算};
    \draw[<-, very thick, warmamber] (-2.2, 0.5) -- (2.2, 0.5);
    \node[below=3pt] at (0, 0.5) {\footnotesize 過程の逆行（Reversing the Process）};
    
    % Bottom quote box safely inside outer rectangle
    \node[draw=darkslate, fill=softgreen, rounded corners=4pt, text width=13.6cm, align=center, inner sep=6pt] at (0, -2.6) {
        \textbf{トンプソンの箴言}：\textit{What one fool can do, another can.}\\
        \footnotesize （一人の人間が理解できる美しい真理は、他のどんな人間にも必ず直観的に理解できる）
    };
\end{tikzpicture}
\caption{微積分の2大基本記号 $d$ と $\int$ の直観的構造と相互関係}
\label{fig:thompson_preliminary}
\end{figure}

教育情報分析研究会　数学教育・物理解析リサーチユニット

\section{【Ollama専任講師特別講義】トンプソン流直観的微積分の扉}

「皆さん、今日から数学の『呪縛』を解き放ちましょう！

多くの生徒が微積分を『複雑な公式の羅列』だと思い込み、身構えています。しかし、シルバヌス・P・トンプソンは第1章でこう断言した。『一人の愚者ができることは、別の愚者にもできる』。これは、微積分が天才の特権ではなく、ただの『論理的な思考の積み重ね』であることを示しています。難解な数式は、単に細分化された思考の集積に過ぎないのです！

では、微積分の核心にある『dx』や『dy』の正体は何でしょうか？ 多くの教科書はこれを曖昧に教えますが、トンプソンの思想に忠実になれば、その正体は『極限まで小さくした、しかしゼロではない変化量』です。

記号 $d$ は「微小な（differential）」という変化の最小単位を指し、$dx$ は「$x$ がほんの少し動いた距離」そのものです。そして積分記号 $\int$ は、ギリシャ文字の $S$（Sum）を長く引き伸ばしたもの。つまり、『無限に小さな変化量を、地道に、しかし正確に積み上げる』という行為の象徴です。

微積分とは、複雑な全体を、扱えるほど小さな断片（$dx$）に分解し、それを再び統合（$\int$）するプロセスです。この『分解と統合』の美学を掴めば、計算の迷宮は一瞬で晴れる。さあ、公式の暗記を捨て、この直観的な真理を掴み取りましょう！」

\section{1. 「一人の馬鹿にできることは、他の馬鹿にもできる」—— トンプソンの挑戦状}

1910年、ロンドン。英国の著名な物理学者であり工学者でもあったシルバヌス・P・トンプソン（Silvanus Phillips Thompson, 1851--1916）は、匿名（F.R.S.＝王立協会フェローの肩書のみ）で一冊の小さな数学書を世に問いました。そのタイトルは『Calculus Made Easy（誰にもわかる微積分）』。そして本の見返しには、古代ローマの古い格言をもじった、痛快で忘れがたい一行が刻まれていました。

「What one fool can do, another can.」

（一人の馬鹿にできることは、他の馬鹿にもできる。）

この言葉こそ、本書の底流に流れるすべての哲学を要約しています。トンプソンは、当時（そして驚くべきことに現代においても）高等教育や高校の現場で微積分が「選ばれた天才しか理解できない秘術」として教えられ、多くの学習者を挫折と絶望に追い込んでいる現状に、激しい憤りを覚えていました。

トンプソンは序文でこう喝破しています。「プロの数学者たちが、自分たちの専門的な領分を守るために、わざと難解な記号や過剰に厳密な抽象概念の城壁を築き、一般の学習者を脅かしているのだ。しかし、微積分の本質的な原理そのものは、驚くほど平易であり、健全な直観を持つ人間なら誰でも完全に理解し、自由に使いこなせる道具なのだ。」

物理学や工学、あるいは今日のデータサイエンスやAIの現場で使われている微積分の核心は、複雑怪奇な数式の迷路ではありません。「物事がほんの少し変化したとき、それに連動して別の物事がどれくらい変化するか」を捉える、極めて人間的で直観的な思考の枠組みなのです。トンプソンは、読者に一切の予備的恐怖を捨て去るよう呼びかけます。どんなに数学が苦手だと自認している人であっても、微積分の扉は広く開かれているのです。

\section{2. 微積分に対する恐怖の正体：過剰な形式主義と難解な記号の壁}

なぜ、多くの高校生や大学生は微積分に対してアレルギーを抱いてしまうのでしょうか。その主たる原因は、19世紀以降に数学界で進められた「極限の厳密化」という歴史的経緯と、教育現場におけるその性急な押し付けにあります。

歴史的に見れば、微積分を創始したアイザック・ニュートン（Isaac Newton, 1642--1727）やゴットフリート・ライプニッツ（Gottfried Wilhelm Leibniz, 1646--1716）は、現代の数学者が求めるような厳密な極限の定義（いわゆる$\varepsilon$-$\delta$論法）を持っていませんでした。彼らは「微小な時間」「微小な増分」という、物理的・幾何学的な直観をフルに活用して、惑星の軌道を計算し、曲線の接線を求め、自然界の法則を解き明かしたのです。厳密な土台が整備されたのは、ニュートンたちから100年以上も後のコーシーやワイエルシュトラスの時代でした。

ところが、学校の授業では、歴史的な発見のワクワク感や直観的なイメージが削ぎ落とされ、最初から「極限値 $\lim_{\Delta x \to 0}$ の計算規則」や抽象的な定義が導入されがちです。分母がゼロに近づくとはどういうことか、なぜゼロで割ってはいけないのにゼロとみなしてよいのか、といった疑問が解消されないまま計算ドリルだけが進むため、学習者は「よくわからない呪文を暗記している」という感覚に陥ってしまいます。

トンプソンは、この「過剰な形式主義の壁」を打ち破るために本書を書きました。彼のアプローチは明快です。まずはニュートンやライプニッツが持っていた素朴で生き生きとした直観に立ち返ること。微小量を恐れず、それを手で触れられる具体的な道具として扱うこと。これこそが、真の数学的理解への最短ルートなのです。

\section{3. 2大基本記号 $d$ と $\int$ の直観的正体：「かけら」と「寄せ集め」}

微積分の教科書を開いたとき、最も威圧感を与えるのが「$d$」と「$\int$」という2つの不可思議な記号です。トンプソンは原著の第1章において、この2つの記号の「正体」を、小学生でもわかる言葉で一気に解き明かしています。

まず、記号「$d$」です。これはラテン語の *differentia*（差、区別）あるいは英語の *difference* の頭文字に由来します。トンプソンはこれを「a little bit of（ほんの少しのかけら）」と読むよう指示します。

\begin{itemize}\item $dx$ とは、「$x$ という量の、ほんの少しのかけら（微小部分）」のことです。\end{itemize}

\begin{itemize}\item $dy$ とは、「$y$ という量の、ほんの少しのかけら」のことです。\end{itemize}

\begin{itemize}\item $dt$ とは、「時間 $t$ の、ほんの少しの瞬間（微小時間）」のことです。\end{itemize}

$dx$ は、$d$ と $x$ の掛け算ではありません。「$x$ の微小な増分」を一文字で表したラベルです。通常の長さが定規で測れるメートル単位だとすれば、$dx$ は顕微鏡でしか見えないほどの、極めて微細な長さのひとかけらを意味します。

次に、記号「$\int$」です。これはアルファベットの「S」を上下に細長く引き伸ばしたものです。なぜ「S」なのかといえば、ラテン語の *summa*（総和、足し算）、英語の *sum* の頭文字だからです。

ライプニッツが考案したこの記号は、文字通り「微小なかけらをすべて集めて足し合わせる」という操作を意味します。

\begin{itemize}\item $\int dx$ とは、「$x$ の微小なかけら $dx$ を、端から端まですべて寄せ集めて足し合わせる」ということであり、結果として元の全体量 $x$ に戻ります。\end{itemize}

つまり、微分（$d$）とは「全体を微小なかけらに切り分ける操作」であり、積分（$\int$）とは「切り分けられた微小なかけらを元通りに足し合わせる操作」に他なりません。この2つは、ちょうど「リンゴを細かくすりおろすこと」と「すりおろしたリンゴを集めて元の重さに戻すこと」のような、互いに完全な逆操作の関係にあるのです。この直観さえ頭に入れば、微積分の呪文のような記号は、親しみやすい道具へと一変します。

\section{4. 高校数学と科学の扉：スピードメーターと変化の瞬間}

この「かけら」と「寄せ集め」という発想が、なぜ現代の数学や物理学、そして科学技術の根幹を支えているのでしょうか。その最も身近な例が、自動車の「スピードメーター」です。

あなたが車に乗って、東京から静岡までの $180\,\mathrm{km}$ を3時間で走ったとします。単純に割り算をすると、

$$\frac{180\,\mathrm{km}}{3\,\text{時間}} = 60\,\mathrm{km/h}$$

となります。これは「平均の速度」です。しかし、道中には赤信号で止まっていた時間もあれば、高速道路を時速 $100\,\mathrm{km}$ で快走していた時間もあるはずです。運転席のスピードメーターが指しているのは、3時間の平均値ではなく、「まさに今、この瞬間の速度」です。

では、「瞬間の速度」とは一体どうやって定義し、測定するのでしょうか。

一瞬とは、時間が全く経過しない「ゼロ秒」のことでしょうか。もし時間がゼロなら、進む距離もゼロですから、$\frac{0}{0}$ となって計算できません。ゼロで割ることは数学では禁止されています。

そこで人類が生み出した大発明が「微小時間 $dt$」と「微小距離 $ds$」です。

1時間ではなく、1分でもなく、1秒、あるいは1000分の1秒という「ほんの少しの時間のかけら $dt$」を考えます。その極めて短い瞬間に、車がどれだけ進んだかという「距離のかけら $ds$」を測ります。$dt$ も $ds$ も極めて微小な量ですが、その比率である

$$\frac{ds}{dt}$$

は、決してゼロにはならず、例えば「時速 $60\,\mathrm{km}$」という、しっかりとした測定可能な有限の数値として定まります。

これこそが、高校数学IIや数学IIIで学ぶ「導関数（変化率）」の正体です。自然界のあらゆる現象――ボールの弾道、惑星の運行、感染症の拡大、株価の変動、スマートフォンの音声処理、AIの学習アルゴリズム――はすべて、「微小な変化 $dx$ に対する、結果の微小な変化 $dy$ の比率 $\frac{dy}{dx}$」として記述されています。微積分とは、止まった世界を眺める静的な数学ではなく、移ろいゆく現実世界そのものを捉えるための「動的な科学の言語」なのです。

\begin{examplebox}{高校生のための重要例題：第1章 —— 微小量 $dx$ と $dy$ の直観的把握}
\textbf{【問題】}\\
時速 $60\,\mathrm{km}$ で等速直線運動している自動車がある。走行時間 $t$（時間）と走行距離 $s$（$\mathrm{km}$）の関係は $s = 60t$ である。いま、微小時間 $dt = 1\,\text{秒} = \frac{1}{3600}\,\text{時間}$ が経過したとき、走行距離の微小増分 $ds$ を求め、比 $\frac{ds}{dt}$ の物理的意味を答えよ。
\\[2mm]
\textbf{【詳細な解答・解説】}\\
関数 $s(t) = 60t$ において、時間 $t$ が $t + dt$ に変化したとき、走行距離は $s + ds = 60(t + dt) = 60t + 60dt$ となる。両辺から $s = 60t$ を引くと、微小増分 $ds = 60dt$ が得られる。$dt = \frac{1}{3600}\,\text{時間}$ を代入すると、$ds = 60 \times \frac{1}{3600} = \frac{1}{60}\,\mathrm{km} \approx 16.67\,\mathrm{m}$ である。このとき、比 $\frac{ds}{dt}$ は $\frac{60dt}{dt} = 60\,\mathrm{km/h}$ となり、微小時間の極限においても一定の「瞬間の速度（スピードメーターの指示値）」を表す。微小量 $ds$ と $dt$ はともに極小であるが、その比 $\frac{ds}{dt}$ は有限で確固たる物理的実在（速度）を示す。
\end{examplebox}


\section*{第1章の確認演習}
\addcontentsline{toc}{section}{第1章の確認演習}
\begin{enumerate}[label=\textbf{問\arabic*.}]
\item 記号 $dx$ と $\Delta x$ の概念的相違を、直観的微小量と有限な増分の観点から説明せよ。
\item トンプソンの箴言「What one fool can do, another can」が示す、微積分学習において直観が形式主義に先行すべき理由を論ぜよ。
\item 積分記号 $\int$ がなぜ文字「S」の引き伸ばしとして考案されたか、その歴史的背景を述べよ。
\end{enumerate}


\chapter{小ささの階級と無視の数理：1次微小量・2次微小量と近似の本質}
\chaptermark{小ささの階級と無視の数理}

\begin{flushright}
\small\textit{On Different Degrees of Smallness}
\end{flushright}

\begin{figure}[htbp]
\centering
\begin{tikzpicture}[x=1cm, y=1cm, >=stealth]
    % Original Square
    \fill[softblue] (0, 0) rectangle (3.5, 3.5);
    \draw[thick, primaryblue] (0, 0) rectangle (3.5, 3.5);
    \node at (1.75, 1.75) {\large\textbf{元の正方形} $x^2$};
    
    % Strips dx
    \fill[softamber] (0, 3.5) rectangle (3.5, 4.3);
    \draw[thick, warmamber] (0, 3.5) rectangle (3.5, 4.3);
    \node at (1.75, 3.9) {\footnotesize 高さ $dx$ の帯：$x \cdot dx$};
    
    \fill[softamber] (3.5, 0) rectangle (4.3, 3.5);
    \draw[thick, warmamber] (3.5, 0) rectangle (4.3, 3.5);
    \node[rotate=-90] at (3.9, 1.75) {\footnotesize 幅 $dx$ の帯：$x \cdot dx$};
    
    % Corner (dx)^2
    \fill[softrose] (3.5, 3.5) rectangle (4.3, 4.3);
    \draw[thick, crimson] (3.5, 3.5) rectangle (4.3, 4.3);
    
    \draw[->, thick, crimson] (4.9, 4.7) -- (4.05, 4.05);
    \node[above right, crimson] at (4.4, 4.6) {\footnotesize 極小正方形 $(dx)^2$（2次微小量）};
    
    % Dimensions
    \draw[<->, thick] (0, -0.4) -- (3.5, -0.4) node[midway, below] {$x$};
    \draw[<->, thick] (3.5, -0.4) -- (4.3, -0.4) node[midway, below] {$dx$};
    \draw[<->, thick] (-0.4, 0) -- (-0.4, 3.5) node[midway, left] {$x$};
    \draw[<->, thick] (-0.4, 3.5) -- (-0.4, 4.3) node[midway, left] {$dx$};
    
    % Expansion Box
    \node[draw=accentgreen, fill=softgreen, rounded corners=4pt, text width=6.2cm, inner sep=7pt] at (7.8, 2.1) {
        \textbf{微小増分の展開}\\[4pt]
        $(x + dx)^2 = x^2 + 2x dx + (dx)^2$\\[4pt]
        \footnotesize $\bullet$ 1次微小量：$2x dx$（本質的増分）\\
        $\bullet$ 2次微小量：$(dx)^2$（桁違いに極小）\\[4pt]
        \normalsize \textbf{結論}：極限 $dx \to 0$ において $(dx)^2$ は\\
        安全に無視され、$\frac{d(x^2)}{dx} = 2x$ となる！
    };
\end{tikzpicture}
\caption{正方形の幾何学的拡大と2次微小量 $(dx)^2$ の自然消滅}
\label{fig:thompson_degrees_smallness}
\end{figure}

教育情報分析研究会　数学教育・物理解析リサーチユニット

\section{【Ollama専任講師特別講義】トンプソン流直観的微積分の扉}

「皆さん、ようこそ！今日は、微積分の『魔法』の正体を暴きに来ました。

多くの人は『微積分は難しい』と怯えますが、シルバヌス・P・トンプソンは断言します。『微積分は、単に、非常に小さなものを扱う技術に過ぎない』と。その核心にあるのが、この第2章の『小ささの階級（Degrees of Smallness）』です。

想像してください。大富豪が、1ファージング（極小単位）の貨幣を落としたとします。それは、彼にとっての『微小なもの』です。しかし、もし彼がさらにその1ファージングを分割した『微小な中の一部』を落としたらどうでしょう？ それはもはや、彼にとっての『無視できるほど小さなもの』となります。これが「階級」です。

これを数式に当てはめてみましょう。正方形の辺を $x$ から $x+dx$ へ広げた時、面積の変化は $(x+dx)^2 = x^2 + 2x dx + (dx)^2$ となります。

ここで重要なのは、$dx$ が「極めて小さい」という性質です。

$dx$ が限りなくゼロに近づくとき、$dx$ は「微小」ですが、その2乗である $(dx)^2$ は、「微小なものよりもさらに圧倒的に小さい」という階級に属します。

例えるなら、$dx$ が『1円』なら、$(dx)^2$ は『1円の100万分の1』のようなものです。極限 $\lim_{dx \to 0}$ を取る際、この「圧倒的な差」により、$(dx)^2$ は主要な変化（$2x dx$）に比べて無視しても良いほど無力な存在になるのです。

微積分とは、この『大きさの階級』を正しく見極め、本質的なものだけを抽出する芸術なのです！」

\section{1. 日常生活における「無視」の知恵：大金持ちと1ファージングの貨幣}

トンプソンは『Calculus Made Easy』の第2章において、微積分を学ぶ上で最も重要でありながら、最も誤解されやすい概念――「微小量の無視（neglecting small quantities）」について、ユーモア溢れる日常のたとえ話を用いて解説しています。

当時の英国の通貨単位には、ポンド（$\pounds$）、シリング（$\mathrm{s}$）、ペンス（$\mathrm{d}$）、そしてペンスのさらに4分の1である「ファージング（farthing, 銅貨）」がありました。1ポンドは20シリング、240ペンス、すなわち960ファージングに相当します。

トンプソンはこう問いかけます。

「もし何百万ポンドもの資産を持つ大富豪が、街を歩いていて道端に1ファージング（数円相当）を落としたとして、それを帳簿にわざわざ記載して資産額を修正するだろうか？ ましてや、そのファージングのさらに1000分の1の微粒子を気にする人がいるだろうか？」

日常生活において、私たちは常に「何が重要で、何を安全に無視できるか」という判断を行っています。例えば、体重計に乗って自分の体重を測るとき、ポケットに微細なホコリが1粒入っているかどうかを気にする人はいません。地球の質量を計算する物理学者が、道端に小石が1個転がったことで地球の重さが変わったと大騒ぎすることもありません。

微積分における「極限の計算」とは、まさにこの健全な「無視の知恵」を、極めて精緻かつ数学的に体系化したものなのです。微積分で切り捨てられる微小量は、「面倒だから適当に丸めた」という雑な近似ではありません。考察している本質的な変化量と比べて「桁違いに極小（高次の微小量）」であるがゆえに、厳密にゼロとみなして差し支えない、必然的な数学的帰結なのです。

\section{2. 微小量の階級（Degrees of Smallness）：1次、2次、3次微小量}

トンプソンは、微小量には明確な「階級（Degrees）」が存在することを強調します。これを理解することが、微積分の計算でなぜ項が消えるのかを納得するための鍵となります。

いま、ある有限の基準となる量（通常の大きさ）を $1$ としましょう。そして、それと比べて極めて微小な量を $dx$ と置きます。

例えば、$dx = \frac{1}{1000} = 10^{-3}$（1000分の1、あるいは0.001）だとイメージしてください。この $dx$ を「第1次の微小量（First order small quantity）」と呼びます。1メートルのものに対する1ミリメートルのかけらです。

では、この微小量どうしを掛け合わせた $(dx)^2 = dx \cdot dx$ はどうなるでしょうか。

$$(dx)^2 = \left(\frac{1}{1000}\right)^2 = \frac{1}{1000000} = 10^{-6}$$

100万分の1（0.000001）になります。これは1メートルのものに対する1マイクロメートル（バクテリアのサイズ）です。

元の微小量 $dx$（ミリメートル）と比べても、さらに1000倍も小さくなっています。これを「第2次の微小量（Second order small quantity）」と呼びます。

さらに、$(dx)^3$ はどうでしょうか。

$$(dx)^3 = \left(\frac{1}{1000}\right)^3 = \frac{1}{1000000000} = 10^{-9}$$

10億分の1（0.000000001、ナノメートル、原子のサイズ）です。これを「第3次の微小量」と呼びます。

ここで重要な原則が導かれます。

\begin{enumerate}\item 有限の量（定数や $x$）に対しては、$dx$ などの1次微小量は極限において無視できます。\end{enumerate}

\begin{enumerate}\item しかし、変化率（比率）を議論する現場では、$dx$ は主役となります。\end{enumerate}

\begin{enumerate}\item そして、1次微小量 $dx$ が主役として活躍する舞台において、2次微小量 $(dx)^2$ や3次微小量 $(dx)^3$ は、「主役の1000分の1、100万分の1」という途方もない小ささになるため、完全に消し去って（無視して）何の影響も及ぼさないのです。\end{enumerate}

\section{3. 正方形の幾何学的展開：$(x+dx)^2 = x^2 + 2x dx + (dx)^2$ の図形的直観}

この「小ささの階級」を、幾何学的に最も美しく視覚化したのが、トンプソンが示した「正方形の拡大モデル」です。

一辺の長さが $x$ である正方形を考えてみましょう。この正方形の面積は、当然ながら

$$S = x^2$$

です。

いま、この正方形の縦と横の長さを、それぞれほんの少しだけ、微小量 $dx$ だけ引き伸ばしたとします。新しい正方形の一辺は $x + dx$ になります。

新しく拡大された正方形の全体の面積は、代数的には次のように展開されます。

$$(x + dx)^2 = x^2 + 2x\,dx + (dx)^2$$

この展開式の各項を、図形として分解して観察してみましょう。

\begin{enumerate}\item 元の正方形：面積 $x^2$（左下の大部分を占める領域）\end{enumerate}

\begin{enumerate}\item 右側に付け足された長方形の帯：幅 $dx$、高さ $x$ なので、面積は $x \cdot dx$\end{enumerate}

\begin{enumerate}\item 上側に付け足された長方形の帯：幅 $x$、高さ $dx$ なので、面積は $x \cdot dx$\end{enumerate}

\begin{enumerate}\item 右上の角に残された極小の正方形：幅 $dx$、高さ $dx$ なので、面積は $(dx)^2$\end{enumerate}

正方形をわずかに拡大したときの「面積の増分」を $\Delta S$ と書くと、

$$\Delta S = (x + dx)^2 - x^2 = 2x\,dx + (dx)^2$$

となります。

面積の増加の大部分は、元の正方形の2つの縁に沿って伸びた細長い2本の帯、すなわち $2x\,dx$（1次の微小量）です。

そして右上の角にある小さな四角形 $(dx)^2$ は、微小量と微小量の積であり、2次の微小量です。

もし $x = 10\,\mathrm{cm}$ で、$dx = 0.01\,\mathrm{cm}$（$0.1\,\mathrm{mm}$）だとすれば、

\begin{itemize}\item 2本の帯の面積 $2x\,dx = 2 \times 10 \times 0.01 = 0.2\,\mathrm{cm}^2$\end{itemize}

\begin{itemize}\item 右上の極小正方形 $(dx)^2 = (0.01)^2 = 0.0001\,\mathrm{cm}^2$\end{itemize}

です。右上の正方形は、帯の面積のわずか2000分の1に過ぎません。

\section{4. 極限 $dx \to 0$ における「消滅」の正当性と近代微積分の厳密性}

さて、ここからが微積分の核心です。

私たちが知りたいのは、「辺の長さ $x$ がほんの少し伸びたとき、面積 $S$ はどれほどの比率（勢い）で増大するか」という変化率 $\frac{dS}{dx}$ です。

増分の式 $\Delta S = 2x\,dx + (dx)^2$ の両辺を、辺の伸び $dx$ で割ってみましょう。

$$\frac{\Delta S}{dx} = \frac{2x\,dx + (dx)^2}{dx} = 2x + dx$$

この式の右辺を見てください。第1項は $2x$ であり、第2項にはまだ $dx$ が残っています。

ここで、伸び $dx$ を限りなくゼロに近づける極限操作（$dx \to 0$）を行います。

すると、残っていた $dx$ は文字通りゼロへと消滅します。

したがって、極限における変化率（導関数）は、

$$\frac{dS}{dx} = 2x$$

と、極めてシンプルかつ厳密に求まるのです。

これが、高校数学で公式として暗記させられる「$(x^2)' = 2x$」の幾何学的証明の全貌です。

初学者が抱く「$(dx)^2$ を途中で勝手に捨ててしまっていいのか？」「それはインチキではないのか？」という疑問に対する答えは明確です。

比率 $\frac{\Delta S}{dx}$ を計算した段階で、$(dx)^2$ は $dx$ で割られてただの $1$ 次の微小量 $dx$ に格下げされます。そして、有限の量 $2x$ の隣に並んだ $dx$ は、$dx \to 0$ の極限において完全に $0$ に収束するため、数学的に寸分の狂いもなく正確に $2x$ だけが残るのです。

トンプソンはこの鮮やかな論理を指して、「微積分とは、高次の微小量を正当に振り落とす技術である」と語りました。難解な哲学に惑わされる必要はありません。大きさに桁違いの差があるとき、極限において小さいものは自然に消え去る――この物理的リアリズムこそが、微積分の強力無比な計算力を支えているのです。

\begin{examplebox}{高校生のための重要例題：第2章 —— 小ささの階級と高次微小量の無視}
\textbf{【問題】}\\
一辺の長さが $x = 10\,\mathrm{cm}$ の正方形の金属板がある。温度上昇により各辺が $dx = 0.01\,\mathrm{cm}$（$0.1\,\mathrm{mm}$）だけ熱膨張した。全体の面積の正確な増加分 $\Delta S$ を求め、1次微小量 $2x dx$ と2次微小量 $(dx)^2$ の数値を比較せよ。
\\[2mm]
\textbf{【詳細な解答・解説】}\\
正方形の面積 $S = x^2$ に対し、膨張後の面積は $(x+dx)^2 = x^2 + 2x dx + (dx)^2$ である。面積の増分は $\Delta S = 2x dx + (dx)^2$ である。数値を代入すると、$2x dx = 2 \times 10 \times 0.01 = 0.2\,\mathrm{cm}^2$ であり、$(dx)^2 = (0.01)^2 = 0.0001\,\mathrm{cm}^2$ である。全体増分 $\Delta S = 0.2001\,\mathrm{cm}^2$ のうち、$(dx)^2$ が占める割合はわずか $\frac{0.0001}{0.2} = 0.05\%$ に過ぎない。$dx \to 0$ の極限では、この割合は完全にゼロへ収束する。したがって、高次の微小量 $(dx)^2$ を消去し、$dS = 2x dx$、すなわち $\frac{dS}{dx} = 2x$ と近似・極限計算を行うことは、完全に合理的かつ数学的に厳密である。
\end{examplebox}


\section*{第2章の確認演習}
\addcontentsline{toc}{section}{第2章の確認演習}
\begin{enumerate}[label=\textbf{問\arabic*.}]
\item 一辺 $x$ の立方体の体積 $(x+dx)^3$ を展開し、なぜ2次および3次の微小量が極限で消去されるか説明せよ。
\item 100万ポンドに対する1ペニーの比率を計算し、微小量の無視が近似としてどれほど高精度か考察せよ。
\item 物理学における「1次近似（微小振動の線形化）」の重要性を述べよ。
\end{enumerate}


\chapter{相対的な成長度：変化率としての導関数 $\frac{dy}{dx}$ と曲線の傾き}
\chaptermark{相対的な成長度}

\begin{flushright}
\small\textit{On Relative Growings and Geometrical Slopes}
\end{flushright}

\begin{figure}[htbp]
\centering
\begin{tikzpicture}[x=1cm, y=1cm, >=stealth]
    % Top Header Box
    \node[draw=warmamber, fill=softamber, rounded corners=4pt, inner sep=5pt] at (4.0, 4.7) {
        \bfseries 相対的成長度：$\lim_{\Delta x \to 0} \frac{\Delta y}{\Delta x} = \frac{dy}{dx}$（瞬間の変化率）
    };
    
    % Axes
    \draw[->, thick] (0, 0) -- (8.2, 0) node[right] {\footnotesize 入力 $x$};
    \draw[->, thick] (0, 0) -- (0, 3.8) node[left] {\footnotesize 出力 $y$};
    
    % Curve y = 0.05 x^2 + 0.2 x + 0.6
    \draw[very thick, primaryblue, domain=0.8:6.5, samples=50] plot (\x, {0.05*\x*\x + 0.2*\x + 0.6});
    \node[above, primaryblue] at (1.8, 2.2) {\bfseries\footnotesize 曲線 $y = f(x)$};
    
    % Points P and Q
    \coordinate (P) at (2.0, 1.2);
    \coordinate (Q) at (5.2, 2.99);
    
    % Secant line (dashed): slope = (2.99-1.2)/3.2 = 0.56
    \draw[thick, dashed, darkslate] (0.8, 0.53) -- (6.6, 3.78);
    \node[right, darkslate] at (6.6, 3.78) {\bfseries\footnotesize 割線の傾き $\frac{\Delta y}{\Delta x}$};
    
    % Tangent line at P: slope = 2*0.05*2 + 0.2 = 0.4
    \draw[very thick, crimson] (0.8, 0.72) -- (4.8, 2.32);
    % Tangent label placed strictly between x=2.0 and x=5.2 underneath tangent line
    \node[below=3pt, crimson] at (3.8, 1.35) {\bfseries\scriptsize 接線の傾き $\frac{dy}{dx} = f'(x)$};
    
    % Points
    \filldraw[crimson] (P) circle (2pt) node[above left=2pt, darkslate] {$P(x, y)$};
    \filldraw[primaryblue] (Q) circle (2pt) node[below right=2pt, darkslate] {$Q(x+\Delta x, y+\Delta y)$};
    
    % Vertical projection dashed lines
    \draw[dotted, thick] (P) -- (2.0, 0) node[below] {$x$};
    \draw[dotted, thick] (Q) -- (5.2, 0) node[below] {$x + \Delta x$};
\end{tikzpicture}
\caption{曲線の割線から接線への極限推移と瞬間変化率（導関数）}
\label{fig:thompson_slopes}
\end{figure}

教育情報分析研究会　数学教育・物理解析リサーチユニット

\section{【Ollama専任講師特別講義】トンプソン流直観的微積分の扉}

「皆さん、今日から微積分の『真の心臓部』に踏み込みます。トンプソンは『微積分は難解な計算ではない。変化の性質を捉えることだ』と説きました。その核心が、この第3章の『相対的な成長度（Relative Growings）』です。

まず、時計の針を思い出してください。長針は短針の12倍の速さで動きます。これは「速度の比」です。トンプソンは、ある量が『他の量に対してどれだけ速く変化しているか』という相対的な比率こそが、微分の本質であると主張します。

ここで重要なのは、グラフ上の「割線の傾き」から「接線の傾き $dy/dx$」への移行です。

通常、私たちは2点間の傾き（割線）を求めますが、トンプソンはこれを『2点間の距離を限りなく縮めていくプロセス』として捉えます。

割線は『平均的な変化』ですが、2点を重ね合わせ、距離 $dx$ を極限までゼロに近づけるとき、それはその瞬間における『純粋な変化の勢い』、すなわち接線の傾きへと変貌します。

$dy/dx$ とは、単なる計算式ではありません。それは『$x$が動くとき、$y$がどれだけ敏感に反応しているか』という動的な比率なのです。この『相対的な成長』の概念を掴んだ時、微積分は初めて血の通った学問へと変わります。準備はいいですか？ 変化の正体を見に行きましょう！」

\section{1. 「絶対的な量」ではなく「相対的な比率」を見る}

トンプソンは第3章の冒頭で、微積分の本質を「On Relative Growings（相対的な成長度について）」という魅力的な言葉で表現しています。私たちが日常で目にする現象を分析するとき、最も重要なのは「ある量が単独でどれだけ増えたか」という絶対的な変化量ではなく、「ひとつの量が変化したとき、それと連動してもうひとつの量がどれだけの割合で変化したか」という相対的な比率です。

例えば、ある木の苗木が1年間に $30\,\mathrm{cm}$ 成長したとします。これは「時間に対する成長」です。

一方で、アクセルペダルを $1\,\mathrm{cm}$ 踏み込んだとき、車の速度が時速 $20\,\mathrm{km}$ 上がったとします。これは「踏み込み量に対する速度の変化」です。

このように、自然界や工学の現象は、常に「原因となる変数（独立変数 $x$）」と「結果として変化する変数（従属変数 $y$）」のペアとして現れます。

微積分が探求するのは、まさにこの2つの変数の間の「相対的な成長のスピード」です。

$x$ がほんの少しだけ $dx$ だけ成長したとき、$y$ はそれに呼応してどれだけ（$dy$ だけ）成長するか。

その2つのかけらの比率である

$$\frac{dy}{dx}$$

を、トンプソンは「相対的成長度（Rate of Growing）」と呼び、現代の数学ではこれを「導関数（Derivative）」あるいは「微分係数」と呼びます。分数のように書かれる $\frac{dy}{dx}$ は、単なる記号の飾りではなく、文字通り「$dy$ 割る $dx$」という比率の直観を今なお色濃く宿しているのです。

\section{2. ロンドンの時計台の針：長針と短針の相対的な速度比}

トンプソンは、この「相対的成長度」を直観的に理解させるために、ロンドンの街頭や教会の時計台にある時計の針を例に挙げます。

時計の文字盤には、短針（時針）と長針（分針）があります。

時間が経過するにつれて、両方の針は回転し、それぞれの先端は円周上を移動していきます。

いま、短針の動いた微小な角度を $dh$、長針の動いた微小な角度を $dm$ としてみましょう。

時計の仕組みとして、長針が文字盤を1周（$360^\circ$）回転する間に、短針は1時間分、すなわち文字盤の12分の1（$30^\circ$）しか回転しません。

したがって、どんな微小な瞬間を切り取っても、長針が動く角度 $dm$ は、短針が動く角度 $dh$ の常に12倍です。

$$dm = 12\,dh \iff \frac{dm}{dh} = 12$$

このとき、比率 $\frac{dm}{dh}$ は $12$ という一定の定数になります。

短針に対する長針の「相対的成長度」は $12$ です。逆に、長針に対する短針の相対的成長度は $\frac{dh}{dm} = \frac{1}{12}$ です。

もし針の動きがこのように常に一定の比率であれば、話は極めて単純です。これはグラフで言えば「傾きが一定の直線（1次関数 $y = ax + b$）」の世界です。

しかし、現実の世界の現象の多くは、直線的ではありません。車は加速し、ボールは重力で曲がり、感染症は爆発的に急増します。すなわち、場所（$x$ の値）によって、相対的成長度 $\frac{dy}{dx}$ そのものが刻一刻と変化していくのです。

この「刻一刻と変化する成長の勢い」を、あらゆる地点で完璧に追跡できるようにすることこそが、微積分の真骨頂なのです。

\section{3. グラフ上の幾何学的意味：割線の極限としての接線の傾き}

では、この相対的成長度 $\frac{dy}{dx}$ を、直交座標系（デカルト座標）のグラフの上で視覚化してみましょう。

関数 $y = f(x)$ のグラフとして、滑らかに曲がった曲線を思い浮かべてください。

曲線上の任意の点 $P(x, y)$ をとります。

ここから、$x$ 座標をほんの少しだけ増やして $x + \Delta x$ とした点を $Q(x + \Delta x, y + \Delta y)$ とします。

点 $P$ と点 $Q$ を結ぶ直線を引くと、これは曲線を横切る「割線（Secant line）」となります。

この割線の傾き（Slope）は、中学生でも知っている「変化の割合」の公式そのものです。

$$\text{割線の傾き} = \frac{\Delta y}{\Delta x} = \frac{f(x + \Delta x) - f(x)}{\Delta x}$$

これは、区間 $[x, x + \Delta x]$ における「平均の傾き」です。直角三角形を描けば、底辺の長さが $\Delta x$、高さが $\Delta y$ であり、底辺に対する高さの比率（正接 $\tan\theta$）を表しています。

いま、点 $P$ を固定したまま、点 $Q$ を曲線に沿って点 $P$ に向かって限りなく近づけていきます。これは $\Delta x \to 0$ とする極限操作です。

点 $Q$ が点 $P$ に限りなく近づくにつれて、割線 $PQ$ はどのような動きを見せるでしょうか。

割線は次第に回転していき、極限において、点 $P$ で曲線に「ぴったりと触れ合う1本の直線」――すなわち接線（Tangent line）へと収束していきます。

このとき、底辺 $\Delta x$ は微小量 $dx$ となり、高さ $\Delta y$ は微小量 $dy$ となります。

そして、接線の傾きこそが、

$$\frac{dy}{dx} = \lim_{\Delta x \to 0} \frac{f(x + \Delta x) - f(x)}{\Delta x}$$

なのです。

「曲線の傾きを測る」とは、曲線を顕微鏡で極限まで拡大して見ることと同じです。どんなに複雑に湾曲した曲線であっても、点 $P$ のごくごく近傍を無限大の倍率で拡大すれば、曲がりは見えなくなり、ほぼ完全に「真っ直ぐな線分」に見えます。その極小線分の傾きこそが $\frac{dy}{dx}$ です。微積分とは、曲線を無数の微小な直線の集まりとして理解する幾何学の顕微鏡なのです。

\section{4. 高校数学IIの平均変化率から瞬間変化率へ：$\Delta x \to 0$ の架け橋}

日本の高校数学カリキュラムにおいて、数学IIの『微分と積分』では、まさにこの概念が次のようなステップで導入されます。

\begin{enumerate}\item 平均変化率: 区間 $[a, b]$ における $\frac{f(b) - f(a)}{b - a}$\end{enumerate}

\begin{enumerate}\item 微分係数: $b \to a$（または $h \to 0$）とした極限値 $f'(a) = \lim_{h \to 0}\frac{f(a+h) - f(a)}{h}$\end{enumerate}

\begin{enumerate}\item 導関数: 特定の点 $a$ だけでなく、任意の変数 $x$ における傾きを与える関数 $f'(x)$\end{enumerate}

トンプソンの記法 $\frac{dy}{dx}$ と、日本の教科書で標準的に使われるラグランジュの記法 $f'(x)$ は、表現は異なりますが全く同じ実体を指しています。

しかし、トンプソンのライプニッツ流記法 $\frac{dy}{dx}$ には、学習上きわめて大きな利点があります。

それは、「何を変数として、何を微分しているのか」がひと目でわかるという点です。

例えば、物理学において物体の位置を $s$、時間を $t$ としたとき、

\begin{itemize}\item 速度は「時間に対する位置の変化率」なので $\frac{ds}{dt}$\end{itemize}

\begin{itemize}\item 加速度は「時間に対する速度の変化率」なので $\frac{dv}{dt}$\end{itemize}

と書かれます。これを単に $s'$ や $v'$ と書くよりも、分母に $dt$ があることで「時間で微分したのだ」という物理的意味が鮮明に浮かび上がります。

さらに、$\frac{dy}{dx}$ という分数形式の記法は、後に学ぶ「合成関数の微分（チェイン・ルール）」や「置換積分」において、あたかも分数の約分のように直観的かつ機械的に計算を進められるという、驚異的な威力を発揮します。

トンプソンが本書で一貫してライプニッツの記法を採用したのは、まさにこの「直観的なわかりやすさ」と「計算の扱いやすさ」を初学者に提供するためだったのです。高校生の皆さんも、難解な記号に臆することなく、「$y$ の微小な伸び割る $x$ の微小な伸び」という素朴なイメージを常に胸に抱いて、次の具体的な計算規則へと進んでいきましょう。

\begin{examplebox}{高校生のための重要例題：第3章 —— 曲線の接線の傾きと瞬間変化率}
\textbf{【問題】}\\
放物線 $y = x^2 - 4x + 5$ 上の点 $P(3, 2)$ における接線の傾きを導関数の定義に基づき求め、その点における接線の方程式を求めよ。
\\[2mm]
\textbf{【詳細な解答・解説】}\\
関数 $f(x) = x^2 - 4x + 5$ に対し、$x$ が $x + dx$ に変化したときの変化量は、\\$f(x+dx) = (x+dx)^2 - 4(x+dx) + 5 = x^2 + 2x dx + (dx)^2 - 4x - 4dx + 5$\\となる。$dy = f(x+dx) - f(x) = (2x - 4)dx + (dx)^2$。両辺を $dx$ で割り、$dx \to 0$ とすると、導関数 $\frac{dy}{dx} = 2x - 4$ が得られる。点 $P(3, 2)$ における接線の傾きは $m = 2(3) - 4 = 2$ である。したがって、求める接線の方程式は $y - 2 = 2(x - 3)$、すなわち $y = 2x - 4$ である。
\end{examplebox}


\section*{第3章の確認演習}
\addcontentsline{toc}{section}{第3章の確認演習}
\begin{enumerate}[label=\textbf{問\arabic*.}]
\item 平均変化率 $\frac{\Delta y}{\Delta x}$ と瞬間変化率 $\frac{dy}{dx}$ の違いを、速度計の動作に例えて説明せよ。
\item 放物線 $y = x^2$ の $x = 2$ における接線の方程式を導関数の定義から導け。
\item 関数が「滑らかである（微分可能である）」とは、幾何学的に何を意味するか述べよ。
\end{enumerate}


\chapter{ベキ関数の微分法：$y = x^n$ の微分公式と直観的証明}
\chaptermark{ベキ関数の微分法}

\begin{flushright}
\small\textit{Simplest Cases: Powers of x and Constants}
\end{flushright}

\begin{figure}[htbp]
\centering
\begin{tikzpicture}[x=1cm, y=1cm, >=stealth]
    % Outer Frame
    \draw[thick, fill=softblue, draw=primaryblue, rounded corners=6pt] (-7.6, -4.2) rectangle (7.6, 4.0);
    \node[draw=primaryblue, fill=white, rounded corners=4pt, inner sep=5pt] at (0, 3.2) {\bfseries ベキ関数 $y = x^n$ の微分公式とその完全体系};
    
    % 3 Columns
    \node[draw=primaryblue, fill=white, rounded corners=4pt, text width=3.8cm, align=center, inner sep=5pt] at (-4.8, 0.9) {
        \textbf{自然数指数 $n \in \mathbb{N}$}\\[3pt]
        \small $y = x^1 \implies y' = 1$\\[2pt]
        $y = x^2 \implies y' = 2x$\\[2pt]
        $y = x^3 \implies y' = 3x^2$\\[3pt]
        \footnotesize 一般則：$y' = n x^{n-1}$
    };
    
    \node[draw=warmamber, fill=softamber, rounded corners=4pt, text width=3.8cm, align=center, inner sep=5pt] at (0, 0.9) {
        \textbf{負の整数（反比例）}\\[3pt]
        \small $y = x^{-1} = \frac{1}{x} \implies y' = -\frac{1}{x^2}$\\[2pt]
        $y = x^{-2} = \frac{1}{x^2} \implies y' = -\frac{2}{x^3}$\\[4pt]
        \footnotesize ベキ乗規則が完全成立！
    };
    
    \node[draw=accentgreen, fill=softgreen, rounded corners=4pt, text width=3.8cm, align=center, inner sep=5pt] at (4.8, 0.9) {
        \textbf{有理数指数（根号）}\\[3pt]
        \small $y = x^{1/2} = \sqrt{x} \implies y' = \frac{1}{2\sqrt{x}}$\\[2pt]
        $y = x^{1/3} = \sqrt[3]{x} \implies y' = \frac{1}{3\sqrt[3]{x^2}}$\\[4pt]
        \footnotesize 実数全体 $n \in \mathbb{R}$ へ拡張
    };
    
    % Bottom General Rule Box
    \node[draw=darkslate, fill=white, rounded corners=4pt, text width=13.6cm, align=center, inner sep=6pt] at (0, -2.6) {
        \textbf{トンプソンの普遍則}：定数倍 $c$ を含む一般形\\[2pt]
        \normalsize $\frac{d}{dx}(c x^n) \;=\; c \cdot n x^{n-1}$ は数学全般の礎石である
    };
\end{tikzpicture}
\caption{ベキ関数の微分公式における指数体系（自然数・負数・分数）の統一性}
\label{fig:thompson_powers}
\end{figure}

教育情報分析研究会　数学教育・物理解析リサーチユニット

\section{【Ollama専任講師特別講義】トンプソン流直観的微積分の扉}

皆さん、よく集まってくれました！今日は、微積分という巨大な山脈の核心、「ベキ関数の微分公式」の正体を暴きます。

多くの教科書は「公式を暗記せよ」と命じますが、トンプソンは違います。彼は「微小な変化の正体」を直視せよと説くのです。

まず、 $x^2$ の微小増分を考えましょう。$(x+h)^2 - x^2$ を展開すると $2xh + h^2$ です。ここで、$h$ を極限まで小さくすると、$h^2$ は無視できるほど小さくなり、残るのは $2xh$。これに $h$ で割れば、変化率は $2x$ となります。

$x^3$ ならどうでしょう？ $(x+h)^3 - x^3$ を展開すれば $3x^2h + 3xh^2 + h^3$。同様に $h^2$ 以降を捨てれば、変化率は $3x^2$ です。

この「展開して、高い次数の項を捨てる」という操作こそが、二項定理の本質です！この論理を突き詰めると、指数の正体が見えてきます。

例えば $1/x$（負の指数）や $\sqrt{x}$（分数指数）も、二項定理を拡張すれば、全く同じ理屈で導かれます。

「指数が何であれ、微小な変化の連鎖は同じ構造（多項式の展開）に従う」。

この一貫性こそが数学の美学です。公式は魔法ではありません。微小な世界を正確に記述するための、極めて自然で必然な帰結なのです！

\section{1. 最も単純なケース：$y = x^2$ と $y = x^3$ の展開と規則性の発見}

トンプソンは第4章において、微積分の具体的な計算技術の第一歩として、数学で最も頻繁に登場する「ベキ関数（Power functions）」の微分法を取り上げます。彼は複雑な抽象論から入るのではなく、最も身近で具体的な $y = x^2$ や $y = x^3$ の観察からスタートします。

まず、$y = x^2$ を復習しましょう。第2章で見たように、$x$ を $x + dx$ に増やすと、

$$y + dy = (x + dx)^2 = x^2 + 2x\,dx + (dx)^2$$

となります。元の $y = x^2$ を引き算すると、増分は

$$dy = 2x\,dx + (dx)^2$$

です。ここで2次の微小量 $(dx)^2$ を無視して $dx$ で割ると、

$$\frac{dy}{dx} = 2x$$

が得られます。指数にあった「$2$」が係数として前に降りてきて、$x$ の次数が1つ減って $x^1 = x$ になっています。

次に、$y = x^3$（立方体の体積など）を調べてみましょう。

一辺が $x$ の立方体の体積は $V = x^3$ です。一辺が $dx$ だけ伸びると、

$$y + dy = (x + dx)^3 = x^3 + 3x^2\,dx + 3x\,(dx)^2 + (dx)^3$$

となります。両辺から $y = x^3$ を引くと、体積の微小増分 $dy$ は

$$dy = 3x^2\,dx + 3x\,(dx)^2 + (dx)^3$$

となります。

幾何学的に見ると、立方体の拡大は次の3つの要素から成り立っています：

\begin{enumerate}\item 3つの面に貼り付く薄い板：それぞれ面積 $x^2$、厚さ $dx$ なので、合計体積は $3x^2\,dx$（1次微小量）\end{enumerate}

\begin{enumerate}\item 3つの角の稜線にできる細長い角材：体積 $3x\,(dx)^2$（2次微小量）\end{enumerate}

\begin{enumerate}\item 頂点の小さなサイコロ：体積 $(dx)^3$（3次微小量）\end{enumerate}

ここで $dx$ で割って相対的成長度を求めると、

$$\frac{dy}{dx} = 3x^2 + 3x\,dx + (dx)^2$$

となります。極限 $dx \to 0$ において、第2項と第3項は $dx$ を含むため綺麗に消滅します。

したがって、残るのは

$$\frac{dy}{dx} = 3x^2$$

です。ここでも、元の指数「$3$」が前に降りてきて、次数が1つ減って $x^2$ になっています。

\section{2. 一般の整数ベキ $y = x^n$：二項定理と1次微小量の係数 $n$}

$x^2$ の微分が $2x$、$x^3$ の微分が $3x^2$ となれば、誰もが「$x^4$ の微分は $4x^3$、一般の $x^n$ の微分は $n x^{n-1}$ になるのではないか？」と予想するでしょう。トンプソンは、この直観的予想が二項定理によって完全に正当化されることを示します。

任意の正の整数 $n$ に対して、$(x + dx)^n$ を二項展開してみましょう。高校数学IIで学ぶ二項定理によれば、

$$(x + dx)^n = x^n + \binom{n}{1} x^{n-1}\,dx + \binom{n}{2} x^{n-2}\,(dx)^2 + \dots + (dx)^n$$

となります。ここで $\binom{n}{1} = n$ です。

したがって、

$$y + dy = x^n + n x^{n-1}\,dx + \frac{n(n-1)}{2} x^{n-2}\,(dx)^2 + \dots$$

元の $y = x^n$ を引き算すると、増分 $dy$ は

$$dy = n x^{n-1}\,dx + \left[ (dx)^2 \text{ 以上の高次微小量の項} \right]$$

となります。

両辺を $dx$ で割ると、

$$\frac{dy}{dx} = n x^{n-1} + \left[ dx \text{ を因数に含む項} \right]$$

となります。ここで極限 $dx \to 0$ をとれば、第2項以降の無数の項は、すべて $dx$ を掛け算されているためにことごとく $0$ へと吸い込まれて消滅します。

かくして、燦然と輝く微積分の基本公式が得られます：

$$\frac{d}{dx}(x^n) = n x^{n-1}$$

トンプソンはこの公式の美しさをこう絶賛しています。「肩に乗っている指数 $n$ を前に引きずり下ろして掛け算し、肩の指数を $1$ つ減らす。これほど簡単で、しかも万能な計算規則が他にあるだろうか！」

\section{3. 負の指数への拡張：$y = x^{-1} = \frac{1}{x}$ の微分と反比例}

この公式 $\frac{d}{dx}(x^n) = n x^{n-1}$ の驚くべきところは、指数 $n$ が正の整数（$1, 2, 3, \dots$）のときだけでなく、負の整数や分数（有理数）、さらには任意の実数に対しても全く同じ形で成立するという点です。

まずは、反比例のグラフでおなじみの関数 $y = \frac{1}{x} = x^{-1}$ を考えてみましょう。

微小量を用いたトンプソン流の素朴な計算を実行してみます。

$x$ を $x + dx$ に増やすと、

$$y + dy = \frac{1}{x + dx}$$

となります。増分 $dy$ は、通分して引き算することで求まります：

$$dy = \frac{1}{x + dx} - \frac{1}{x} = \frac{x - (x + dx)}{x(x + dx)} = \frac{-dx}{x^2 + x\,dx}$$

両辺を $dx$ で割ると、

$$\frac{dy}{dx} = \frac{-1}{x^2 + x\,dx}$$

となります。ここで極限 $dx \to 0$ をとると、分母の $x\,dx$ は消滅し、$x^2$ だけが残ります。

したがって、

$$\frac{dy}{dx} = -\frac{1}{x^2} = -1 \cdot x^{-2}$$

が得られます。

この結果を、ベキ関数の公式 $\frac{d}{dx}(x^n) = n x^{n-1}$ と見比べてみてください。

$n = -1$ を代入すると、

$$\frac{d}{dx}(x^{-1}) = (-1) x^{-1 - 1} = -x^{-2} = -\frac{1}{x^2}$$

となり、完全に一致しています。

負の符号が出る理由も直観的に納得できます。反比例のグラフ $y = 1/x$ は、$x$ が増えるにつれて $y$ の値が減少する（右肩下がりの）曲線です。したがって、成長度（傾き）はマイナスでなければならないのです。

同様に、$y = x^{-2} = \frac{1}{x^2}$ であれば、微分すると $-2x^{-3} = -\frac{2}{x^3}$ となります。指数が負であっても、公式は完璧に機能します。

\section{4. 分数指数（有理数）への拡張：$y = \sqrt{x} = x^{1/2}$ と根号の微分}

次に、平方根（ルート）を含む関数 $y = \sqrt{x} = x^{1/2}$ について考えます。高校生がつまずきやすい根号の微分も、トンプソンの手にかかれば実にエレガントに解明されます。

$y = \sqrt{x}$ の両辺を2乗すると、

$$y^2 = x$$

となります。ここで両辺の微小量を考えてみましょう。

左辺の $y^2$ を微分すると、$2y\,dy$ になります（$y$ の正方形の面積増分）。

右辺の $x$ の微小量は $dx$ です。

したがって、

$$2y\,dy = dx$$

という極めてシンプルな関係が成り立ちます。

私たちが求めたいのは相対的成長度 $\frac{dy}{dx}$ ですから、両辺を変形すると、

$$\frac{dy}{dx} = \frac{1}{2y}$$

となります。ここで元の定義 $y = \sqrt{x}$ を代入して戻せば、

$$\frac{dy}{dx} = \frac{1}{2\sqrt{x}} = \frac{1}{2} x^{-1/2}$$

が得られます。

これをベキ関数の公式で直接計算してみましょう。

$n = \frac{1}{2}$ として公式を適用すると、

$$\frac{d}{dx}\left(x^{1/2}\right) = \frac{1}{2} x^{\frac{1}{2} - 1} = \frac{1}{2} x^{-1/2} = \frac{1}{2\sqrt{x}}$$

となり、全く同じ結果が一瞬で導かれます。

一般に、3乗根 $y = \sqrt[3]{x} = x^{1/3}$ であっても、$y = x^{p/q}$ であっても、

$$\frac{d}{dx}(x^r) = r x^{r-1} \quad (r \in \mathbb{R})$$

という単一の統一公式が支配しています。

高校数学IIIでは、この公式を丸暗記して機械的に計算させられることが多いですが、その背景には「両辺の微小増分の釣り合い」という極めて自然で直観的な幾何学的・代数学的構造が横たわっているのです。

\begin{examplebox}{高校生のための重要例題：第4章 —— ベキ関数の微分公式の応用}
\textbf{【問題】}\\
次の各関数の導関数を求めよ。\\(1) $y = 5x^4 - 2x^3 + 7x - 9$ \quad (2) $y = \frac{3}{x^2}$ \quad (3) $y = 4\sqrt{x}$
\\[2mm]
\textbf{【詳細な解答・解説】}\\
(1) ベキ関数の基本公式 $\frac{d}{dx}(x^n) = n x^{n-1}$ および定数項の微分はゼロである規則より、\\$\frac{dy}{dx} = 5(4x^3) - 2(3x^2) + 7(1) - 0 = 20x^3 - 6x^2 + 7$。\\(2) 負の指数表示 $y = 3x^{-2}$ に直す。公式を適用して、\\$\frac{dy}{dx} = 3(-2x^{-3}) = -6x^{-3} = -\frac{6}{x^3}$。\\(3) 分数指数表示 $y = 4x^{1/2}$ に直す。公式より、\\$\frac{dy}{dx} = 4\left(\frac{1}{2}x^{-1/2}\right) = 2x^{-1/2} = \frac{2}{\sqrt{x}}$。
\end{examplebox}


\section*{第4章の確認演習}
\addcontentsline{toc}{section}{第4章の確認演習}
\begin{enumerate}[label=\textbf{問\arabic*.}]
\item 二項定理を用いて、自然数 $n$ に対する公式 $\frac{d}{dx}(x^n) = n x^{n-1}$ を証明せよ。
\item 関数 $y = \frac{1}{x^3}$ および $y = \sqrt[3]{x}$ の導関数を求めよ。
\item 定数関数の微分がゼロになる理由を、幾何学的傾きの観点から説明せよ。
\end{enumerate}


\chapter{関数の結合と演算規則：和・差・積・商の微分公式のメカニズム}
\chaptermark{関数の結合と演算規則}

\begin{flushright}
\small\textit{Sums, Differences, Products, and Quotients}
\end{flushright}

\begin{figure}[htbp]
\centering
\begin{tikzpicture}[x=1cm, y=1cm, >=stealth]
    % Original Rectangle u x v
    \fill[softblue] (0, 0) rectangle (3.8, 2.8);
    \draw[thick, primaryblue] (0, 0) rectangle (3.8, 2.8);
    \node at (1.9, 1.4) {\large\textbf{元の面積} $y = u \cdot v$};
    
    % Top strip: u * dv
    \fill[softamber] (0, 2.8) rectangle (3.8, 3.6);
    \draw[thick, warmamber] (0, 2.8) rectangle (3.8, 3.6);
    \node at (1.9, 3.2) {\footnotesize 上の増加帯：$u \cdot dv$};
    
    % Right strip: v * du (width 1.4cm for comfortable text)
    \fill[softamber] (3.8, 0) rectangle (5.2, 2.8);
    \draw[thick, warmamber] (3.8, 0) rectangle (5.2, 2.8);
    \node[align=center] at (4.5, 1.4) {\footnotesize 右の帯\\[2pt] $v \cdot du$};
    
    % Corner: du * dv
    \fill[softrose] (3.8, 2.8) rectangle (5.2, 3.6);
    \draw[thick, crimson] (3.8, 2.8) rectangle (5.2, 3.6);
    
    % Label for corner
    \draw[->, thick, crimson] (4.5, 4.5) -- (4.5, 3.7);
    \node[above, crimson] at (4.5, 4.5) {\footnotesize 極小長方形 $du \cdot dv$（2次微小量）};
    
    % Dimensions
    \draw[<->, thick] (0, -0.4) -- (3.8, -0.4) node[midway, below] {$u$};
    \draw[<->, thick] (3.8, -0.4) -- (5.2, -0.4) node[midway, below] {$du$};
    \draw[<->, thick] (-0.4, 0) -- (-0.4, 2.8) node[midway, left] {$v$};
    \draw[<->, thick] (-0.4, 2.8) -- (-0.4, 3.6) node[midway, left] {$dv$};
    
    % Algebra Box at Right (starting cleanly at x=6.8)
    \node[draw=accentgreen, fill=softgreen, rounded corners=4pt, text width=6.2cm, inner sep=6pt] at (10.0, 2.0) {
        \textbf{積の増分公式}\\[3pt]
        $d(uv) = (u + du)(v + dv) - uv$\\[2pt]
        $\phantom{d(uv)} = u \cdot dv + v \cdot du + du \cdot dv$\\[3pt]
        \footnotesize $\bullet$ 1次増分：$u \cdot dv + v \cdot du$\\
        $\bullet$ 2次増分：$du \cdot dv \to 0$（消去）\\[4pt]
        \normalsize \textbf{両辺を $dx$ で割ると：}\\[3pt]
        $\frac{d(uv)}{dx} = u \frac{dv}{dx} + v \frac{du}{dx}$
    };
\end{tikzpicture}
\caption{長方形の面積変化による積の微分公式 $(uv)' = u'v + uv'$ の幾何学的証明}
\label{fig:thompson_product_rule}
\end{figure}

教育情報分析研究会　数学教育・物理解析リサーチユニット

\section{【Ollama専任講師特別講義】トンプソン流直観的微積分の扉}

「皆さん、注目してください！今日は微積分の『魔法』の正体を暴きます。多くの人が公式を丸暗記して終わらせますが、トンプソンは『なぜそうなるのか』という動的な変化を捉えることを求めています。

まず積の微分 $ (uv)' = u'v + uv' $ です。

想像してください。縦の長さ $u$、横の長さ $v$ の長方形があります。この両方を同時に、ごくわずかに（微小量 $\Delta u, \Delta v$）伸ばしたとき、面積の変化はどうなるでしょう？

面積の増分は、『縦を伸ばしたことによる増分（$\Delta u \cdot v$）』と『横を伸ばしたことによる増分（$u \cdot \Delta v$）』の合計です。これを微分の言葉に変換すれば、まさに $u'v + uv'$。面積の増加は、二つの方向への伸びの積み重ねなのです。

では、商の微分 $\left(\frac{u}{v}\right)'$ になぜマイナスが現れるのか。

ここが本質です。分母 $v$ を大きくすることは、全体の値を「小さくする」こと、つまり負の変化を意味します。

商の微小変化を計算すると、分子の変化 $u \cdot \Delta v$ と分母の変化 $\Delta u \cdot v$ が現れますが、分母の変化は「分母を大きくして全体を減らす」という逆の作用を働かせます。この『分母の増加による減少』という逆転現象が、公式におけるマイナス符号の正体です。

公式は単なる記号の羅列ではありません。面積の伸びや、分母による抑制といった物理的な動きを、数学が正確に記述しているのです！」

\section{1. 和と差の微分：線形性と足し算の独立性}

単独のベキ関数 $x^n$ の微分法をマスターした私たちが次に直面するのは、複数の関数が足し算、引き算、掛け算、割り算によって結合された複合的な数式の微分です。トンプソンは第5章で、これらの演算規則を「独立した変化の足し合わせ」という物理的観点から解き明かします。

最も直観的で平易なのが「和と差の微分」です。

いま、ある量 $y$ が、2つの独立した部分 $u$ と $v$ の和として表されているとします。

$$y = u + v$$

ここで、$u$ も $v$ もそれぞれ変数 $x$ に依存して変化する関数だとします。

$x$ が微小量 $dx$ だけ変化したとき、$u$ は微小量 $du$ だけ変化し、$v$ は微小量 $dv$ だけ変化します。

その結果、全体の量 $y$ は次のように変化します：

$$y + dy = (u + du) + (v + dv) = (u + v) + (du + dv)$$

両辺から $y = u + v$ を引くと、全体の微小増分 $dy$ は単純に2つの増分の足し算になります：

$$dy = du + dv$$

両辺を $dx$ で割れば、直ちに次の公式が得られます：

$$\frac{dy}{dx} = \frac{du}{dx} + \frac{dv}{dx}$$

引き算 $y = u - v$ の場合も全く同様で、$\frac{dy}{dx} = \frac{du}{dx} - \frac{dv}{dx}$ となります。

また、定数倍 $y = c \cdot u$（$c$ は定数）であれば、$dy = c \cdot du$ ですから、$\frac{dy}{dx} = c \frac{du}{dx}$ です。

これは数学でいう「線形性（Linearity）」と呼ばれる性質です。物理的に言えば、2つのバネが直列に繋がっているとき、全体の伸びはそれぞれのバネの伸びの単純な和になる、という当たり前の事実を表しています。多項式 $y = 3x^4 - 5x^2 + 7x - 2$ のような関数を微分するとき、各項をバラバラに微分して足し引きしてよいのは、この線形性が保証されているからに他なりません。

\section{2. 積の微分公式 $(uv)' = u'v + uv'$：長方形の面積拡大モデルによる直観的証明}

足し算が単純な和になるのに対し、多くの学習者が最初の壁として突き当たるのが「積の微分（Product Rule）」です。

「$(u \cdot v)'$ は $u' \cdot v'$ になるのではないか？」という素朴な誤解は、歴史上の大数学者ライプニッツ自身も初期に犯した有名な間違いです。しかし、少し立ち止まって考えれば、それが誤りであることは明白です。

トンプソンは、この積の微分のメカニズムを、第2章の正方形を一般化した「長方形の面積拡大モデル」によって鮮やかに視覚化します。

いま、横の長さが $u$、縦の長さが $v$ である長方形を考えます。この長方形の面積は、

$$y = u \cdot v$$

です。

変数 $x$ が微小量 $dx$ だけ増えたとき、横の長さは微小量 $du$ だけ伸び、縦の長さは微小量 $dv$ だけ伸びるとします。

拡大された新しい長方形の面積は、次のようになります：

$$y + dy = (u + du)(v + dv) = uv + v\,du + u\,dv + du\,dv$$

両辺から元の面積 $y = uv$ を差し引くと、面積の正味の増分 $dy$ は次の3つの領域の和として表されます：

$$dy = v\,du + u\,dv + du\,dv$$

これを図形で確認してみましょう：

\begin{enumerate}\item 右側に付け足された細長い長方形：幅 $du$、高さ $v$ なので、面積は $v\,du$\end{enumerate}

\begin{enumerate}\item 上側に付け足された細長い長方形：幅 $u$、高さ $dv$ なので、面積は $u\,dv$\end{enumerate}

\begin{enumerate}\item 右上の隅にできる極小の長方形：幅 $du$、高さ $dv$ なので、面積は $du \cdot dv$\end{enumerate}

\section{3. 2次微小量 $du \cdot dv$ の自然な脱落と美しい対称性}

ここで、第2章で学んだ「小ささの階級（Degrees of Smallness）」の原理が威力を発揮します。

$du$ も $dv$ も、それぞれ $dx$ に比例する「1次の微小量」です。

したがって、右上の隅にある極小長方形 $du \cdot dv$ は、微小量どうしの積、すなわち「2次の微小量」となります。

相対的成長度を求めるために、増分の式全体を $dx$ で割ってみましょう：

$$\frac{dy}{dx} = v \frac{du}{dx} + u \frac{dv}{dx} + du \frac{dv}{dx}$$

第3項を見てください。$\frac{dv}{dx}$ は有限の変化率ですが、その前に微小量 $du$ が掛け算されています。

ここで極限 $dx \to 0$（すなわち $du \to 0$）をとると、第3項は完全にゼロへと消滅します！

こうして、後に残るのは第1項と第2項だけとなります：

$$\frac{d(uv)}{dx} = u \frac{dv}{dx} + v \frac{du}{dx}$$

あるいは、プライム記号を用いて書けば、

$$(uv)' = u'v + uv'$$

となります。

なんと美しい対称性でしょう！

「積の微分は、片方を微分してもう片方はそのまま、逆を微分してもう片方はそのまま、そしてそれらを足し合わせる。」

この公式の形を丸暗記するのではなく、「横が伸びたことによる面積の増加（$v\,du$）」と「縦が伸びたことによる面積の増加（$u\,dv$）」という2枚の短冊の足し合わせとして捉えることで、公式は一生忘れられない生きた知識となるのです。

\section{4. 商の微分公式 $(u/v)' = (u'v - uv')/v^2$：分母の増大がもたらすマイナスの効果}

積の微分を理解した後に待ち受けるのが、分数の形をした関数、すなわち「商の微分（Quotient Rule）」です。

$$y = \frac{u}{v}$$

高校数学で最も符号や分母を間違えやすいこの公式も、トンプソンの論理に従えば明快に導出できます。

まず、素朴な微小量の変化から出発してみましょう。

$$y + dy = \frac{u + du}{v + dv}$$

増分 $dy$ を計算するために通分して引き算を行います：

$$dy = \frac{u + du}{v + dv} - \frac{u}{v} = \frac{v(u + du) - u(v + dv)}{v(v + dv)} = \frac{vu + v\,du - uv - u\,dv}{v^2 + v\,dv} = \frac{v\,du - u\,dv}{v^2 + v\,dv}$$

両辺を $dx$ で割ると、

$$\frac{dy}{dx} = \frac{v \frac{du}{dx} - u \frac{dv}{dx}}{v^2 + v\,dv}$$

となります。分母にある $v\,dv$ は微小量ですから、極限 $dx \to 0$（$dv \to 0$）において消滅し、分母は単なる $v^2$ になります。

したがって、商の微分公式が得られます：

$$\frac{d}{dx}\left(\frac{u}{v}\right) = \frac{u'v - uv'}{v^2}$$

ここで極めて重要な直観的ポイントは、なぜ分子にマイナス符号（$-uv'$）が現れるのかという点です。

分数 $\frac{u}{v}$ において、分子 $u$ が増えれば全体の値は増えます（プラスの寄与、$v\,u'$）。

しかし、分母 $v$ が増えれば、割り算の性質上、全体の値は減少します（マイナスの寄与、$-uv'$）。

ケーキを分けるとき、ケーキの量（分子）が増えれば取り分は増えますが、食べる人数（分母）が増えれば1人あたりの取り分は減ってしまうのと同じです。

この「分母の成長がもたらす逆方向の引っ張り」が、分子のマイナス符号となって数式に刻み込まれているのです。数式の形には、必ず物理的・日常的な必然性が宿っています。

\begin{examplebox}{高校生のための重要例題：第5章 —— 積の微分法と商の微分法の実践}
\textbf{【問題】}\\
次の各関数を微分せよ。\\(1) $y = (x^2 + 3)(2x - 1)$ \quad (2) $y = \frac{x^2 - 1}{x^2 + 1}$
\\[2mm]
\textbf{【詳細な解答・解説】}\\
(1) 積の微分公式 $(uv)' = u'v + uv'$ を用いる。$u = x^2 + 3$, $v = 2x - 1$ とおくと、$u' = 2x$, $v' = 2$。\\$\frac{dy}{dx} = (2x)(2x - 1) + (x^2 + 3)(2) = 4x^2 - 2x + 2x^2 + 6 = 6x^2 - 2x + 6$。\\（※展開してから微分した $y = 2x^3 - x^2 + 6x - 3 \implies y' = 6x^2 - 2x + 6$ と完全に一致する）。\\(2) 商の微分公式 $\left(\frac{u}{v}\right)' = \frac{u'v - uv'}{v^2}$ を用いる。$u = x^2 - 1$, $v = x^2 + 1$ より $u' = 2x$, $v' = 2x$。\\$\frac{dy}{dx} = \frac{2x(x^2 + 1) - (x^2 - 1)(2x)}{(x^2 + 1)^2} = \frac{(2x^3 + 2x) - (2x^3 - 2x)}{(x^2 + 1)^2} = \frac{4x}{(x^2 + 1)^2}$。
\end{examplebox}


\section*{第5章の確認演習}
\addcontentsline{toc}{section}{第5章の確認演習}
\begin{enumerate}[label=\textbf{問\arabic*.}]
\item 長方形の面積モデルを用いて、積の微分公式 $(uv)' = u'v + uv'$ を幾何学的に図解・説明せよ。
\item 商の微分公式を用いて、関数 $y = \frac{2x+1}{x-3}$ を微分せよ。
\item 3つの関数の積 $(uvw)'$ の公式を導出せよ。
\end{enumerate}


\chapter{合成関数と連鎖律：「便利な裏技」チェイン・ルールの威力}
\chaptermark{合成関数と連鎖律}

\begin{flushright}
\small\textit{Introducing a Useful Dodge: The Chain Rule}
\end{flushright}

\begin{figure}[htbp]
\centering
\begin{tikzpicture}[x=1cm, y=1cm, >=stealth]
    % Outer Frame
    \draw[thick, fill=softamber!30, draw=warmamber, rounded corners=6pt] (-7.6, -4.2) rectangle (7.6, 4.0);
    \node[draw=warmamber, fill=white, rounded corners=4pt, inner sep=5pt] at (0, 3.2) {\bfseries 合成関数の微分法（連鎖律・チェインルール）のメカニズム};
    
    % 3 Variable Nodes
    \node[draw=primaryblue, fill=white, rounded corners=4pt, text width=3.4cm, align=center, inner sep=6pt] at (-5.0, 0.9) {
        \textbf{入力変数 $x$}\\[3pt]
        \footnotesize 微小変動 $dx$\\[4pt]
        \scriptsize すべての変化の源泉
    };
    
    \node[draw=accentgreen, fill=softgreen, rounded corners=4pt, text width=3.4cm, align=center, inner sep=6pt] at (0, 0.9) {
        \textbf{中間変数 $u = g(x)$}\\[3pt]
        \footnotesize 伝播量 $du = g'(x)dx$\\[4pt]
        \scriptsize ギア1の伝達比 $\frac{du}{dx}$
    };
    
    \node[draw=crimson, fill=softrose, rounded corners=4pt, text width=3.4cm, align=center, inner sep=6pt] at (5.0, 0.9) {
        \textbf{出力変数 $y = f(u)$}\\[3pt]
        \footnotesize 最終変動 $dy = f'(u)du$\\[4pt]
        \scriptsize ギア2の伝達比 $\frac{dy}{du}$
    };
    
    % Interconnecting Arrows with clear above/below labels
    \draw[->, very thick, warmamber] (-3.1, 0.9) -- (-1.9, 0.9);
    \node[above=4pt] at (-2.5, 0.9) {\bfseries\footnotesize $\frac{du}{dx}$};
    
    \draw[->, very thick, warmamber] (1.9, 0.9) -- (3.1, 0.9);
    \node[above=4pt] at (2.5, 0.9) {\bfseries\footnotesize $\frac{dy}{du}$};
    
    % Bottom Formula Box
    \node[draw=darkslate, fill=white, rounded corners=4pt, text width=13.6cm, align=center, inner sep=6pt] at (0, -2.6) {
        \textbf{連鎖律の代数公式}：\Large $\frac{dy}{dx} \;=\; \frac{dy}{du} \cdot \frac{du}{dx}$\\[3pt]
        \footnotesize （中間変数 $u$ の微小変化率を掛け合わせることで、合成関数の微分は一瞬で解ける）
    };
\end{tikzpicture}
\caption{中間変数 $u$ を介した変化率の積としての連鎖律（チェイン・ルール）}
\label{fig:thompson_chain_rule}
\end{figure}

教育情報分析研究会　数学教育・物理解析リサーチユニット

\section{【Ollama専任講師特別講義】トンプソン流直観的微積分の扉}

皆さん、注目してください！今日は微積分の「真の魔法」の一つ、連鎖律（チェイン・ルール）を、トンプソンの思想に基づいて解剖します。

多くの学生は、複雑な関数の微分を「公式の暗記」として捉え、そこに絶望します。しかし、トンプソンの『Calculus Made Easy』が教えてくれる本質は違います。「複雑なものは、単純なものの組み合わせに過ぎない」という視点です。

例えば、$y = f(g(x))$ という関数を考えてみましょう。これは、一見すると巨大で手に負えない塊に見えます。しかし、ここで「中間変数 $u$」という概念を導入してください。$u = g(x)$ と置くのです。すると、この問題は「$x$ から $u$ への変化」と「$u$ から $y$ への変化」という、2つの単純なステップに分解されます。

微分とは「変化の割合」です。$x$ が動いた時、まず $u$ がどれだけ動くか（$du/dx$）。そして、その動いた $u$ を受けて $y$ がどれだけ動くか（$dy/du$）。この2つの変化率を掛け合わせるだけで、複雑な関数の全容が暴かれるのです。

$dy/dx = (dy/du) \cdot (du/dx$

これは単なる計算テクニックではありません。「複雑な関数の皮を、中間変数というメスを使って一枚ずつ剥ぎ取っていく」という、数学的思考の極致です。この「分解の美学」を理解した瞬間、微積分はただの計算から、世界の構造を解き明かす強力な武器へと変貌するのです！

\section{1. トンプソンが明かす「極めて便利な裏技（A Useful Dodge）」}

トンプソンは『Calculus Made Easy』の第6章において、次のような魅力的な告白で読者を引き込みます。

「数学者たちは、難解な理論を駆使して問題を解いているように見せかけるのが好きだが、実際には彼ら自身も、日常的にある『極めて便利な裏技（A Useful Dodge）』を使って複雑な計算を一瞬で片付けている。その裏技の名を『合成関数の連鎖律（Chain Rule）』という。」

高校数学IIIを学ぶ学習者にとって、「合成関数の微分」は最初の大きな関門とみなされがちです。$(3x^2 - 5)^4$ や $\sqrt{x^2 + 1}$ のような式を前にして、多くの生徒は「4乗を展開しなければいけないのか？」「根号をどう扱えばいいのか？」と頭を抱えてしまいます。

しかし、トンプソンが教える「連鎖律」の精神を身につければ、どんなに複雑に見える数式であっても、まるでマトリョーシカ人形の外側の皮を一枚ずつ剥いていくように、極めて機械的かつ爽快に微分できるようになります。

\section{2. 中間変数 $u$ の導入と分数の約分のような連鎖 $\frac{dy}{dx} = \frac{dy}{du} \cdot \frac{du}{dx}$}

連鎖律の幾何学的・代数学的メカニズムは、驚くほどシンプルです。

いま、3つの量 $x, u, y$ があり、$x$ の変化が $u$ を動かし、その $u$ の変化が $y$ を動かすという「玉突き事故」のような連鎖関係を考えてみましょう。

\begin{itemize}\item $x$ がほんの少し $dx$ だけ動いたとき、$u$ は $du$ だけ動く。その比率は $\frac{du}{dx}$ である。\end{itemize}

\begin{itemize}\item $u$ がほんの少し $du$ だけ動いたとき、$y$ は $dy$ だけ動く。その比率は $\frac{dy}{du}$ である。\end{itemize}

では、最初の一押しである $dx$ に対する、最終的な結果 $dy$ の比率 $\frac{dy}{dx}$ はどうなるでしょうか。

微小量のかけらの比率として見れば、当たり前の等式が成り立ちます：

$$\frac{dy}{dx} = \frac{dy}{du} \cdot \frac{du}{dx}$$

右辺を見てください。中間変数のかけらである $du$ が、あたかも分数の掛け算のように分母と分子で「約分」されて消え、左辺の $\frac{dy}{dx}$ と一致することが直観的に見て取れます。

厳密な現代数学においては、「$\frac{dy}{dx}$ は分数ではなく極限の記号であるから、単純に約分してはならない」という注意がなされます。しかし、微小増分の比率の極限として厳密に証明したとしても、得られる結果はこの「あたかも約分できるかのような形」と寸分の狂いもなく一致します。

トンプソンはこれこそがライプニッツの天才的な記法設計の賜物であると強調します。「意味を正しく理解している限り、学習者はこの記法の便利さを最大限に活用して、分数を掛けるように計算して何ら恥じることはない！」

\section{3. 高校数学IIIの最重要関門：複雑な多項式ベキ $(ax+b)^n$ や $(x^2+1)^3$ の一撃攻略}

この「裏技」を、実際の問題に適用してみましょう。

例えば、高校数学で頻出する次の関数を微分してみます：

$$y = (3x^2 - 5)^4$$

もし連鎖律を知らなければ、括弧の中身を4乗展開して膨大な手計算をしなければなりません。しかし、連鎖律を使えば暗算レベルで解くことができます。

括弧の中身という「塊（かたまり）」を、新しい文字（中間変数）$u$ と置きます。

$$u = 3x^2 - 5$$

すると、元の関数は極めてシンプルなベキ関数に変身します：

$$y = u^4$$

それぞれの微分を計算します：

\begin{enumerate}\item $y$ を $u$ で微分する：$\frac{dy}{du} = 4u^3$（外側の皮の微分）\end{enumerate}

\begin{enumerate}\item $u$ を $x$ で微分する：$\frac{du}{dx} = 6x$（内側の身の微分）\end{enumerate}

あとは連鎖律を使って両者を掛け算するだけです：

$$\frac{dy}{dx} = \frac{dy}{du} \cdot \frac{du}{dx} = (4u^3) \cdot (6x) = 24x\,u^3$$

最後に、$u$ を元の $3x^2 - 5$ に戻せば完了です：

$$\frac{dy}{dx} = 24x(3x^2 - 5)^3$$

わずか数行で、しかも展開のミスを犯すリスクなく完璧な答えが得られました。

この操作は「外側の関数を微分して、中身の微分をポンと外に掛け算する」という合言葉で親しまれています。根号のついた関数 $y = \sqrt{x^2 + 1} = (x^2 + 1)^{1/2}$ であっても、$u = x^2 + 1$ と置けば、

$$\frac{dy}{dx} = \frac{1}{2}(u)^{-1/2} \cdot \frac{du}{dx} = \frac{1}{2\sqrt{x^2 + 1}} \cdot 2x = \frac{x}{\sqrt{x^2 + 1}}$$

と、瞬時に導出できます。

\section{4. 多段連鎖と物理学への応用：位置・温度・時間の三者関係}

連鎖律の威力は、中間変数が1つの場合に留まりません。3つでも4つでも、鎖（チェーン）のようにどこまでも繋げることができます。

$$\frac{dy}{dx} = \frac{dy}{du} \cdot \frac{du}{dv} \cdot \frac{dv}{dx}$$

この多段連鎖は、物理学や工学の現場で日常的に活用されています。

例えば、気球に乗って上空へ上昇していく観測者を考えてみましょう。

\begin{itemize}\item 時間を $t$ とする。\end{itemize}

\begin{itemize}\item 気球の高度を $h$ とする。上昇速度は $\frac{dh}{dt}$ である。\end{itemize}

\begin{itemize}\item 上空の気温を $T$ とする。高度に伴う気温の低下率は $\frac{dT}{dh}$ である。\end{itemize}

地上にいる管制官が「時間の経過に伴って、気球の周囲の温度は毎分何度のペースで冷え込んでいるか（$\frac{dT}{dt}$）」を知りたいとき、直接的な温度と時間の関数が手元になくても、連鎖律を用いれば一撃で求まります：

$$\frac{dT}{dt} = \frac{dT}{dh} \cdot \frac{dh}{dt}$$

すなわち、「高度 $1\,\mathrm{m}$ 上がるごとの温度変化」に「$1$ 分間に気球が上昇する高度」を掛け算すればよいのです。

このように、連鎖律とは単なる計算テクニックではなく、「異なる物理量どうしがどのように因果関係の鎖で結ばれているか」を解き明かすための、極めて普遍的な数理的ツールなのです。トンプソンが「便利な裏技」と呼んだこの手法は、現代のニューラルネットワークにおける誤差逆伝播法（Backpropagation）の数学的基礎そのものであり、今日のAI技術を根底で支える最強のエンジンとなっています。

\begin{examplebox}{高校生のための重要例題：第6章 —— 連鎖律（チェイン・ルール）による合成関数の微分}
\textbf{【問題】}\\
関数 $y = (3x^2 - 5)^4$ の導関数 $\frac{dy}{dx}$ を、中間変数 $u$ を用いた連鎖律によって求めよ。
\\[2mm]
\textbf{【詳細な解答・解説】}\\
中間変数として括弧の中身 $u = 3x^2 - 5$ を設定する。これにより、$y = u^4$ と表される。\\それぞれの微分を計算すると、\\$\frac{dy}{du} = 4u^3, \quad \frac{du}{dx} = 6x$\\連鎖律 $\frac{dy}{dx} = \frac{dy}{du} \cdot \frac{du}{dx}$ を適用すると、\\$\frac{dy}{dx} = 4u^3 \cdot 6x = 24x u^3$\\最後に $u = 3x^2 - 5$ を代入して戻すと、\\$\frac{dy}{dx} = 24x(3x^2 - 5)^3$\\展開することなく、合成関数の微分公式によって極めて簡潔に導出できる。
\end{examplebox}


\section*{第6章の確認演習}
\addcontentsline{toc}{section}{第6章の確認演習}
\begin{enumerate}[label=\textbf{問\arabic*.}]
\item 連鎖律 $\frac{dy}{dx} = \frac{dy}{du}\frac{du}{dx}$ を、ギア比（歯車の連動）のアナロジーを用いて説明せよ。
\item 関数 $y = (2x^3 - 1)^5$ および $y = \sqrt{4 - x^2}$ を微分せよ。
\item 逆関数の微分公式 $\frac{dx}{dy} = \frac{1}{dy/dx}$ の幾何学的意味を論ぜよ。
\end{enumerate}


\chapter{高階微分と運動の物理：第2次導関数 $\frac{d^2y}{dx^2}$、速度・加速度の数理}
\chaptermark{高階微分と運動の物理}

\begin{flushright}
\small\textit{Successive Differentiation and When Time Varies}
\end{flushright}

\begin{figure}[htbp]
\centering
\begin{tikzpicture}[x=1cm, y=1cm, >=stealth]
    % Outer Frame
    \draw[thick, fill=softblue, draw=primaryblue, rounded corners=6pt] (-7.6, -4.2) rectangle (7.6, 4.0);
    \node[draw=primaryblue, fill=white, rounded corners=4pt, inner sep=5pt] at (0, 3.2) {\bfseries 時間変化と第2次導関数：運動力学の微積分階層};
    
    % 3 Kinematic Boxes
    \node[draw=primaryblue, fill=white, rounded corners=4pt, text width=3.4cm, align=center, inner sep=6pt] at (-5.0, 0.9) {
        \textbf{位置 $s(t)$}\\[3pt]
        \footnotesize 空間座標 $[\mathrm{m}]$\\[4pt]
        \scriptsize 自由落下 $s = \frac{1}{2}gt^2$
    };
    
    \node[draw=warmamber, fill=softamber, rounded corners=4pt, text width=3.4cm, align=center, inner sep=6pt] at (0, 0.9) {
        \textbf{速度 $v(t) = \frac{ds}{dt}$}\\[3pt]
        \footnotesize 1次導関数 $[\mathrm{m/s}]$\\[4pt]
        \scriptsize 瞬間の進行ペース $gt$
    };
    
    \node[draw=crimson, fill=softrose, rounded corners=4pt, text width=3.4cm, align=center, inner sep=6pt] at (5.0, 0.9) {
        \textbf{加速度 $a(t) = \frac{d^2s}{dt^2}$}\\[3pt]
        \footnotesize 2次導関数 $[\mathrm{m/s^2}]$\\[4pt]
        \scriptsize 重力加速度 $g$（定数）
    };
    
    % Interconnecting Arrows with clear labels
    \draw[->, very thick, warmamber] (-3.1, 0.9) -- (-1.9, 0.9);
    \node[above=4pt] at (-2.5, 0.9) {\footnotesize 微分 $\frac{d}{dt}$};
    
    \draw[->, very thick, warmamber] (1.9, 0.9) -- (3.1, 0.9);
    \node[above=4pt] at (2.5, 0.9) {\footnotesize 微分 $\frac{d}{dt}$};
    
    % Bottom Newton Box
    \node[draw=darkslate, fill=white, rounded corners=4pt, text width=13.6cm, align=center, inner sep=6pt] at (0, -2.6) {
        \textbf{ニュートンの運動方程式}：\Large $F \;=\; m \cdot a \;=\; m \frac{d^2s}{dt^2}$\\[3pt]
        \footnotesize （力は質量 $\times$ 位置の第2次導関数である）
    };
\end{tikzpicture}
\caption{物体の運動における位置・速度・加速度の微積分的階層構造}
\label{fig:thompson_successive_diff}
\end{figure}

教育情報分析研究会　数学教育・物理解析リサーチユニット

\section{【Ollama専任講師特別講義】トンプソン流直観的微積分の扉}

皆さん、ようこそ！今日は、数学の歴史を塗り替えた傑作『Calculus Made Easy』の核心に触れましょう。

第7章でトンプソンが突きつけるのは、「導関数は単なる記号の操作ではなく、『変化の速度』を捉えるための極めて直感的な道具である」という真理です。

まず、なぜ $\frac{d^2y}{dx^2}$ と書くのか。これは単なる「2乗」ではありません。微分 $\frac{dy}{dx}$ は「$x$ が変化したとき、$y$ がどれだけ変化するか」という「変化率」です。その「変化率自体が、さらに $x$ の変化に対してどう変化するか」を求めるのが2階微分です。つまり、「変化の勢いの変化」を測っているのです。

これを物理学に当てはめると、魔法のように鮮明になります。

位置 $y$ を時間 $x$ で微分すれば、速度 $v = \frac{dy}{dx}$ です。さらにその速度を時間で微分すれば、加速度 $a = \frac{d^2y}{dx^2}$ となります。

ガリレオの自由落下を思い出してください。物体は一定の加速度で加速します。つまり、加速度 $\frac{d^2y}{dx^2}$ が「定数」である状態です。この「2階微分が一定」という数理から、積分（逆演算）を繰り返すことで、速度は「時間に比例する直線」になり、位置は「時間に比例する2次関数（放物線）」になることが導かれます。

微分とは、静止した数式に「動き」を吹き込む行為です。この2階微分の概念を理解した時、皆さんは初めて、宇宙の動的な法則を数式で支配する術を手に入れるのです！

\section{1. 微分を繰り返す：導関数の導関数とは何か}

微積分を学び始めた学習者が次に足を踏み入れる魅力的な世界が、「微分を2回、3回と繰り返す」という高階微分（Successive Differentiation）の領域です。トンプソンは第7章において、この概念を抽象的な数式のお遊びとしてではなく、物理的な「運動（動き）」を記述する不可欠な言語として解説しています。

私たちが関数 $y = f(x)$ を微分して得た導関数 $\frac{dy}{dx}$ は、グラフの各地点における「傾き」や「成長の速さ」を表すものでした。しかし、その傾き自身もまた、$x$ の値が変わるにつれて変化していく新しい関数です。

山道を登っている場面を想像してください。最初は緩やかな上り坂（傾きが小さい）だったものが、進むにつれて急勾配になり（傾きが大きくなる）、頂上に近づくと再び緩やかになっていきます。

このように、「傾きそのものが、どれほどの勢いで変化しているか」を知りたいとき、私たちは導関数 $\frac{dy}{dx}$ をもう一度微分します。これが「第2次導関数（Second Derivative）」です。

\section{2. 記号 $\frac{d^2y}{dx^2}$ の読み方と構造：なぜ $d$ に2がつき、$x$ に2がつくのか}

ここで多くの高校生が抱く素朴な疑問があります。

「なぜ第2次導関数の記号は $\frac{d^2y}{dx^2}$ と書くのか？ 分子は $d$ の右肩に $2$ があるのに、分母は $x$ の右肩に $2$ があるのはなぜなのか？」

一見すると不揃いで奇妙に見えるこの記法も、トンプソンの解説を読めば、極めて論理的で厳密な理由があることが一瞬で氷解します。

第2次導関数とは、記号的に言えば「$\frac{d}{dx}$（$x$ で微分せよという命令・作用素）」を、導関数 $\frac{dy}{dx}$ に対してもう一度作用させることです：

$$\frac{d}{dx}\left(\frac{dy}{dx}\right)$$

この式の分子と分母を、掛け算のように形式的にまとめてみましょう。

\begin{itemize}\item 分子は、「$d$」という微分操作が「$dy$」に対して行われるので、$d$ が2回重なって $d(dy) = d^2 y$ と書かれます。$y$ は従属変数そのものであり、$y$ を2乗しているわけではないので、$2$ は微分記号 $d$ につきます。\end{itemize}

\begin{itemize}\item 一方、分母は、「$dx$」という微小な増分が2回掛け合わされるので、$(dx) \times (dx) = (dx)^2$、括弧を省略して $dx^2$ と書かれます。\end{itemize}

つまり、$\frac{d^2y}{dx^2}$ は「$x$ の微小量の2乗 $(dx)^2$ に対する、$y$ の第2次微小増分 $d^2y$ の比率」を寸分の狂いもなく正確に表しているのです。トンプソンはこの記法の明快さを称賛し、声に出して読むときは「ディー・ツー・ワイ、バイ、ディー・エックス・スクエアード（$d^2y$ by $dx^2$）」とリズミカルに読むようアドバイスしています。

\section{3. 物理学の基本法則：位置 $x(t)$、速度 $v(t) = \dot{x}$、加速度 $a(t) = \ddot{x}$}

この第2次導関数が最も鮮やかにその真価を発揮するのが、高校物理における「物体の運動の解析」です。

時間 $t$ を独立変数として、直線上を運動する物体の位置を $x(t)$ としましょう。

\begin{enumerate}\item 位置 $x(t)$: 物体が今どこにいるかを表す。\end{enumerate}

\begin{enumerate}\item 速度 $v(t) = \frac{dx}{dt}$: 位置の時間に対する第1次導関数。物体の位置が時間とともにどれだけ変化しているか、すなわち「移動の速さと向き（スピードメーターの指示値）」を表す。\end{enumerate}

\begin{enumerate}\item 加速度 $a(t) = \frac{dv}{dt} = \frac{d^2x}{dt^2}$: 速度の時間に対する第1次導関数、すなわち位置の第2次導関数。速度が時間とともにどれだけ増減しているか、すなわち「加速の激しさ（背中がシートに押し付けられる感覚）」を表す。\end{enumerate}

ニュートンは時間微分の記法として、文字の上にドットを打つ「ドット記法」を導入しました。

$$v = \dot{x}, \quad a = \ddot{x}$$

あなたが車に乗っていて、アクセルを一定に踏み込んでいるとき、体感しているのは速度そのものではなく「加速度（第2次導関数）」です。車輪がどれほど高速で回っていても、等速運動（$\frac{dx}{dt} = \text{一定}$）であれば加速度はゼロ（$\frac{d^2x}{dt^2} = 0$）であり、体には何の力も感じません。体が力を感じるのは、速度が変化した瞬間、すなわち第2次導関数が非ゼロになった瞬間だけなのです。

\section{4. ガリレオの自由落下とニュートンの運動方程式 $F = ma = m \frac{d^2x}{dt^2}$}

この関係の最も美しい具体例が、ガリレオ・ガリレイが見出した「自由落下運動」です。

高い塔の上から小球を静かに落としたとき、落下し始めてからの時間 $t$ と、落下した距離 $y(t)$ の関係は、高校物理で学ぶように次式で与えられます（$g$ は重力加速度 $\approx 9.8\,\mathrm{m/s^2}$）：

$$y(t) = \frac{1}{2}gt^2$$

これを時間 $t$ で1回微分してみましょう：

$$v(t) = \frac{dy}{dt} = \frac{1}{2}g \cdot (2t) = gt$$

小球の落下速度は時間に比例してどこまでも速くなります。

では、これをもう一度時間 $t$ で微分してみましょう：

$$a(t) = \frac{d^2y}{dt^2} = \frac{d}{dt}(gt) = g$$

なんと、時間 $t$ が完全に消え去り、一定の定数 $g$ だけが残りました！

これは「どんなに時間が経過しても、小球に働く重力による加速度は常に一定（等加速度）である」という物理的真理を完璧に証明しています。

アイザック・ニュートンが確立した近代科学最大の金字塔、運動方程式：

$$F = ma = m \frac{d^2x}{dt^2}$$

は、文字通り「力 $F$ とは、質量 $m$ に位置の第2次導関数を掛け合わせたものである」という関係を宣言しています。惑星の公転軌道からロケットの打ち上げ、ビルや橋の耐震設計に至るまで、人類の近代文明のあらゆる工学計算は、この第2次導関数を解くこと（微分方程式）によって成り立っているのです。

\begin{examplebox}{高校生のための重要例題：第7章 —— 高階微分と自由落下運動の解析}
\textbf{【問題】}\\
地上から初速度 $v_0 = 19.6\,\mathrm{m/s}$ で鉛直上向きに投げ上げられた小球の時刻 $t$（秒）における高さが $h(t) = 19.6t - 4.9t^2$（$\mathrm{m}$）で与えられている。\\(1) 時刻 $t$ における小球の速度 $v(t) = \frac{dh}{dt}$ と加速度 $a(t) = \frac{d^2h}{dt^2}$ を求めよ。\\(2) 小球が最高点に達する時刻と、その高さを求めよ。
\\[2mm]
\textbf{【詳細な解答・解説】}\\
(1) 速度は高さの1次導関数である：\\$v(t) = \frac{dh}{dt} = 19.6 - 9.8t\,\mathrm{[m/s]}$\\加速度は速度の1次導関数（高さの2次導関数）である：\\$a(t) = \frac{d^2h}{dt^2} = -9.8\,\mathrm{[m/s^2]}$\\これは下向きの一定な重力加速度 $g = 9.8\,\mathrm{m/s^2}$ を示す。\\(2) 最高点では瞬間の速度がゼロになるため、$v(t) = 19.6 - 9.8t = 0$ より $t = 2.0\,\text{秒}$。\\そのときの高さは $h(2) = 19.6(2) - 4.9(4) = 39.2 - 19.6 = 19.6\,\mathrm{m}$ である。
\end{examplebox}


\section*{第7章の確認演習}
\addcontentsline{toc}{section}{第7章の確認演習}
\begin{enumerate}[label=\textbf{問\arabic*.}]
\item 第2次導関数 $\frac{d^2y}{dx^2}$ が、曲線の「凹凸（曲がり具合）」を示す理由を説明せよ。
\item 物体の位置が $s(t) = t^3 - 6t^2 + 9t$ で表されるとき、速度がゼロになる時刻と、そのときの加速度を求めよ。
\item ニュートンの運動方程式 $F = m\frac{d^2x}{dt^2}$ がなぜ「第2階微分方程式」であるか述べよ。
\end{enumerate}


\chapter{極大・極小と最適化：導関数がゼロになる点と現実世界の最大化・最小化}
\chaptermark{極大・極小と最適化}

\begin{flushright}
\small\textit{Maxima and Minima: The Art of Optimization}
\end{flushright}

\begin{figure}[htbp]
\centering
\begin{tikzpicture}[x=1cm, y=1cm, >=stealth]
    % Axes
    \draw[->, thick] (0, 0) -- (8.5, 0) node[right] {$x$};
    \draw[->, thick] (0, 0) -- (0, 4.6) node[above] {$y = f(x)$};
    
    % Cubic curve: Max at x=2.8, Min at x=5.8, Inflection at x=4.3
    \draw[very thick, primaryblue, domain=1.4:6.3, samples=50] 
        plot (\x, {0.12*(\x-4.3)*(\x-4.3)*(\x-4.3) - 0.81*(\x-4.3) + 2.8});
    
    % Local Max: x=2.8, y=3.61
    \coordinate (MaxPt) at (2.8, 3.61);
    \filldraw[crimson] (MaxPt) circle (2pt);
    \draw[thick, dashed, crimson] (1.8, 3.61) -- (3.8, 3.61);
    \draw[dotted, thick] (MaxPt) -- (2.8, 0) node[below] {$x_{\mathrm{max}}$};
    \node[above=8pt, crimson] at (MaxPt) {\bfseries\small 極大値 $\left(\frac{dy}{dx} = 0,\ \frac{d^2y}{dx^2} < 0\right)$};
    
    % Inflection point: x=4.3, y=2.8, callout to top-right empty area
    \coordinate (InfPt) at (4.3, 2.8);
    \filldraw[darkslate] (InfPt) circle (2pt);
    \draw[->, thick, darkslate] (5.4, 3.5) -- (4.45, 2.9);
    \node[above right, darkslate] at (5.3, 3.4) {\footnotesize\textbf{変曲点} $\left(\frac{d^2y}{dx^2} = 0\right)$};
    
    % Local Min: x=5.8, y=1.99
    \coordinate (MinPt) at (5.8, 1.99);
    \filldraw[accentgreen] (MinPt) circle (2pt);
    \draw[thick, dashed, accentgreen] (4.8, 1.99) -- (6.8, 1.99);
    \draw[dotted, thick] (MinPt) -- (5.8, 0) node[below] {$x_{\mathrm{min}}$};
    \node[below=8pt, accentgreen] at (MinPt) {\bfseries\small 極小値 $\left(\frac{dy}{dx} = 0,\ \frac{d^2y}{dx^2} > 0\right)$};
\end{tikzpicture}
\caption{極大・極小における接線の傾き（ゼロ）と2次導関数による凹凸判定}
\label{fig:thompson_maxima_minima}
\end{figure}

教育情報分析研究会　数学教育・物理解析リサーチユニット

\section{【Ollama専任講師特別講義】トンプソン流直観的微積分の扉}

皆さん、ようこそ。今日は数学の「心臓部」に触れに来ました。

シルバヌス・P・トンプソンは『Calculus Made Easy』において、微積分を「複雑な計算の技法」ではなく「変化の理（ことわり）」として捉えました。第8章の極大・極小、すなわち「最適化」こそが、この理の真髄です。

まず、フェルマーの原理を思い出してください。山の頂上や谷の底に立ったとき、足元はどう感じられますか？ わずかな瞬間、地面は「水平」です。数学的に言えば、変化率（微分係数）がゼロ。$dy/dx = 0$ という状態は、関数が「一瞬だけ休み、次への方向を迷っている」極点を示します。これが最適化の第一歩、候補の選別です。

しかし、水平な場所すべてが頂上とは限りません。ここで第2次導関数が登場します。関数の「曲がり方（凹凸）」を測るのです。$d^2y/dx^2$ が負なら、グラフは上に凸（山）であり、正なら下に凸（谷）です。この「曲率」の判定こそが、極大と極小を峻別する決定的な証拠となります。

最後に、この理論をブリキ板の箱作りに適用しましょう。板の四隅を切り取る幅を $x$ とすると、容積 $V$ は $x$ の関数となります。$V$ を $x$ で微分し、それがゼロになる地点を探す。そこが「容積が最大になる魔法の数値」です。

微積分とは、複雑な世界から「最高の瞬間」を抽出する知的な武器なのです。この美しさを、ぜひ皆さんの手で計算によって体感してください！

\section{1. 頂上と谷底の特徴：平らな接線と $\frac{dy}{dx} = 0$}

微積分を学ぶ最大の喜びと実用的な利益のひとつが、「関数の最大値や最小値を完璧に特定する技術」――すなわち最適化（Optimization）の力を手に入れることです。トンプソンは第8章において、この「極大（Maxima）と極小（Minima）」の幾何学的・物理的本質を、山登りの直観を通じて解き明かします。

曲がりくねった山道を登り、山の頂上に立った瞬間を思い浮かべてみてください。

登っている最中は、道は常に上り坂（傾き $\frac{dy}{dx} > 0$）です。

頂上を通り過ぎて反対側に下り始めると、道は下り坂（傾き $\frac{dy}{dx} < 0$）になります。

では、まさに「その山の頂上（最高点）」に立っている足元は、どうなっているでしょうか。

登りでも下りでもない、完全に水平で平らな足場（接線の傾きがゼロ）になっているはずです！

谷底を通過するときも全く同じです。

下り坂から上り坂に切り替わる一番底の地点では、傾きは一瞬だけ水平になります。

グラフの上で言えば、曲線が滑らかに山（極大点）を描くときも、谷（極小点）を描くときも、その頂点における接線は地面と平行、すなわち

$$\frac{dy}{dx} = 0$$

となります。

この「極大・極小では導関数がゼロになる」という原理こそ、17世紀フランスの天才ピエール・ド・フェルマー（Pierre de Fermat, 1607--1665）が発見した、微積分学における最も基本的で強力な定理なのです。

\section{2. フェルマーの原理と極値の判定法：第1次導関数の符号変化}

ある関数 $y = f(x)$ の最大や最小を探したいとき、私たちは闇雲にいろんな $x$ を代入して試す必要はありません。導関数を計算し、方程式

$$\frac{dy}{dx} = 0$$

を解くだけでよいのです。この方程式の解となる点（停留点）の候補の中に、真の極値が必ず潜んでいます。

しかし、注意深い高校生の皆さんはここで気づくはずです。

「傾きがゼロ（$\frac{dy}{dx} = 0$）だからといって、必ずしも山や谷とは限らないのではないか？」

その通りです！ 例えば、$y = x^3$ のグラフを考えてみましょう。

導関数は $\frac{dy}{dx} = 3x^2$ ですから、$x = 0$ で傾きはゼロになります。

しかし、$x = 0$ の前後を調べると、$x < 0$ でも傾きはプラス（$3x^2 > 0$）、$x > 0$ でも傾きはプラスです。グラフは一瞬だけ踊り場で平らになった後、再び上り坂を登っていきます。これは極値ではなく「変曲点（Inflection Point）」と呼ばれる点です。

したがって、真の極大・極小であるかを確定するためには、停留点の前後で傾き $\frac{dy}{dx}$ の符号がどう変化しているかを確かめる必要があります：

\begin{enumerate}\item 極大点（山）: 傾きが「プラス（登り） $\to$ ゼロ $\to$ マイナス（下り）」と変化する。\end{enumerate}

\begin{enumerate}\item 極小点（谷）: 傾きが「マイナス（下り） $\to$ ゼロ $\to$ プラス（登り）」と変化する。\end{enumerate}

\begin{enumerate}\item 極値でない点: 傾きの符号が変わらない（プラス $\to$ ゼロ $\to$ プラス、またはマイナス $\to$ ゼロ $\to$ マイナス）。\end{enumerate}

\section{3. 第2次導関数による山と谷の峻別：$\frac{d^2y}{dx^2} < 0$（山）と $\frac{d^2y}{dx^2} > 0$（谷）}

この極大と極小の判定を、増減表を書く手間を省いて一撃で行うための優雅な手法があります。それこそが、第7章で学んだ「第2次導関数 $\frac{d^2y}{dx^2}$」の活用です。

もう一度、山の頂上付近での「傾き」の変化に注目してみましょう。

頂上を通過するとき、傾きは「正の値」から「ゼロ」を経て「負の値」へと、減少し続けています。

傾きという関数が減少しているということは、その導関数（傾きの変化率＝第2次導関数）はマイナス（負）でなければなりません！

逆に、谷底を通過するときは、傾きは「負の値」から「ゼロ」を経て「正の値」へと、増加し続けています。したがって、傾きの変化率（第2次導関数）はプラス（正）になります。

これより、極めて明快な判定法が得られます：

\begin{itemize}\item $\frac{dy}{dx} = 0$ かつ $\frac{d^2y}{dx^2} < 0$ のとき $\implies$ 極大（山）（上に凸）\end{itemize}

\begin{itemize}\item $\frac{dy}{dx} = 0$ かつ $\frac{d^2y}{dx^2} > 0$ のとき $\implies$ 極小（谷）（下に凸）\end{itemize}

第2次導関数が負であることは、曲線が「雨傘のように上に凸」であることを意味し、正であることは「お椀のように下に凸」であることを意味します。符号と形状のイメージを結びつければ、どちらが極大でどちらが極小かを絶対に混同することはありません。

\section{4. トンプソンの古典的名作問題：四隅を切り取って作る箱の容積最大化}

トンプソンは『Calculus Made Easy』の中で、この最適化理論の威力を実感させるために、工学・ものづくりの現場で実際に直面する数多くの名作問題を提示しています。その代表格が「正方形の金属板から作る蓋なし箱の容積最大化問題」です。

【問題】一辺の長さが $a$ の正方形のブリキ板がある。四隅から合同な小さな正方形（一辺の長さ $x$）を切り落とし、残った四方を折り曲げて、蓋のない直方体の箱を作る。箱の容積が最大になるような切り取り幅 $x$ を求めよ。

【トンプソン流の解法】

箱の底面は一辺の長さが $a - 2x$ の正方形となり、箱の高さは $x$ になります。

したがって、箱の容積 $V$ は $x$ の関数として次のように表されます：

$$V(x) = x(a - 2x)^2 = x(a^2 - 4ax + 4x^2) = 4x^3 - 4ax^2 + a^2 x$$

切り取れる幅 $x$ の物理的な範囲は、$0 < x < \frac{a}{2}$ です。

容積を最大にするため、$V$ を $x$ で微分してゼロとおきます：

$$\frac{dV}{dx} = 12x^2 - 8ax + a^2 = 0$$

この二次方程式を因数分解すると、

$$(2x - a)(6x - a) = 0$$

解は $x = \frac{a}{2}$ または $x = \frac{a}{6}$ です。

$x = \frac{a}{2}$ では板が全てなくなってしまうので、求める解の候補は $x = \frac{a}{6}$ です。

本当に極大（最大）であるかを第2次導関数で検証しましょう：

$$\frac{d^2V}{dx^2} = 24x - 8a$$

$x = \frac{a}{6}$ を代入すると、

$$\frac{d^2V}{dx^2} = 24\left(\frac{a}{6}\right) - 8a = 4a - 8a = -4a < 0$$

見事に負（マイナス）となり、上に凸、すなわち極大値であることが確定しました！

したがって、板の一辺のちょうど6分の1を切り取ったとき、箱の容積は最大になるのです。

勘や試行錯誤に頼ることなく、たった数行の微分の計算によって、「材料を最も無駄なく使い、最大の容量を得るための黄金比」が完璧に導かれました。これこそが、数学が現実世界に提供する最強の知性なのです。

\begin{examplebox}{高校生のための重要例題：第8章 —— 容積最大化問題（トンプソンの名作古典問）}
\textbf{【問題】}\\
一辺の長さが $12\,\mathrm{cm}$ の正方形のブリキ板がある。四隅から合同な正方形（一辺 $x\,\mathrm{cm}$）を切り取り、折り曲げて蓋のない直方体の箱を作る。箱の容積 $V$ が最大となるときの切り取り幅 $x$ と、その最大容積 $V_{\text{max}}$ を求めよ。
\\[2mm]
\textbf{【詳細な解答・解説】}\\
箱の底面は一辺 $(12 - 2x)\,\mathrm{cm}$ の正方形、高さは $x\,\mathrm{cm}$ である。ただし $0 < x < 6$。\\容積関数は $V(x) = x(12 - 2x)^2 = x(144 - 48x + 4x^2) = 4x^3 - 48x^2 + 144x$。\\極値を求めるため導関数を計算する：\\$\frac{dV}{dx} = 12x^2 - 96x + 144 = 12(x^2 - 8x + 12) = 12(x - 2)(x - 6)$\\$0 < x < 6$ において停留点は $x = 2$ のみ。\\第2次導関数を調べると、$\frac{d^2V}{dx^2} = 24x - 96$ より、$x = 2$ で $\frac{d^2V}{dx^2} = 48 - 96 = -48 < 0$（上に凸、極大）。\\したがって $x = 2\,\mathrm{cm}$ のとき容積は最大となり、最大容積は $V(2) = 2(12 - 4)^2 = 2 \times 64 = 128\,\mathrm{cm}^3$ である。
\end{examplebox}


\section*{第8章の確認演習}
\addcontentsline{toc}{section}{第8章の確認演習}
\begin{enumerate}[label=\textbf{問\arabic*.}]
\item 極大点と極小点において、なぜ必ず1次導関数がゼロ（$\frac{dy}{dx} = 0$）になるかフェルマーの原理から説明せよ。
\item 関数 $y = x^3 - 3x$ の増減表を作成し、極大値・極小値を求めよ。
\item 2次導関数を用いた極値の判別法（$\frac{d^2y}{dx^2} \gtrless 0$）の有効性と限界を論ぜよ。
\end{enumerate}


\chapter{自然の成長と指数・対数：ネイピア数 $e$、真の複利法則と自然対数の微分}
\chaptermark{自然の成長と指数・対数}

\begin{flushright}
\small\textit{On True Compound Interest and the Exponential Magic}
\end{flushright}

\begin{figure}[htbp]
\centering
\begin{tikzpicture}[x=1cm, y=1cm, >=stealth]
    % Axes
    \draw[->, thick] (-3.8, 0) -- (4.0, 0) node[right] {$x$};
    \draw[->, thick] (0, 0) -- (0, 4.8) node[above] {$y$};
    
    % Curve y = e^x
    \draw[very thick, primaryblue, domain=-3.5:1.35, samples=60] plot (\x, {exp(\x)});
    \node[right, primaryblue] at (1.35, 3.8) {\bfseries 指数関数 $y = e^x$};
    
    % Point at x=1, y=e=2.718
    \coordinate (EPt) at (1.0, 2.718);
    \filldraw[crimson] (EPt) circle (2pt);
    
    % Tangent line: y = e*x
    \draw[thick, dashed, crimson] (0, 0) -- (1.4, 3.805);
    \node[right, crimson] at (1.4, 2.6) {\footnotesize 接線（傾き $= e$）};
    
    % Projections with y=e clearly above dashed line
    \draw[dotted, thick] (EPt) -- (1.0, 0);
    \draw[dotted, thick] (EPt) -- (0, 2.718);
    \node[above=2pt, primaryblue] at (0.5, 2.718) {\footnotesize $y = e$};
    
    % Box at left (fully clear of y axis)
    \node[draw=accentgreen, fill=softgreen, rounded corners=4pt, text width=2.9cm, align=center, inner sep=6pt] at (-2.3, 3.4) {
        \textbf{自己増殖の奇跡}\\[4pt]
        $\frac{d}{dx}(e^x) = e^x$\\[4pt]
        \footnotesize 傾き $=$ 高さ
    };
    
    % Subtangent at bottom
    \node[below=2pt] at (1.0, 0) {$x = 1$};
    \draw[<->, thick, warmamber] (0, -0.6) -- (1.0, -0.6);
    \node[below=3pt, warmamber] at (0.5, -0.6) {\footnotesize 次接線長 $= 1$（一定）};
\end{tikzpicture}
\caption{ネイピア数 $e$ の指数関数 $y = e^x$ における「傾き$=$高さ」の幾何学的自己同一性}
\label{fig:thompson_exp_magic}
\end{figure}

教育情報分析研究会　数学教育・物理解析リサーチユニット

\section{【Ollama専任講師特別講義】トンプソン流直観的微積分の扉}

皆さん、注目してください！今日は数学の歴史において最も美しい「必然」の一つを、トンプソン流の視点で解き明かします。

まず、皆さんに「複利」を想像してください。銀行の利息のように、利息にさらに利息がつく仕組みです。もし、利息を1年ごとではなく、1ヶ月、1日、1秒……と、無限に細かく（連続的に）積み上げ続けたら、お金はどこまで増えるのか？

ここで登場するのが、ネイピア数 $e$ です。複利の回数を無限大に飛ばした極限の到達点、それが $e \approx 2.71828$ です。これは単なる定数ではありません。「変化のスピードが、その瞬間の大きさそのものに等しい」という、自然界の成長ルールを司る基礎定数なのです。

この $e$ を使った関数 $f(x) = e^x$ は、数学界の奇跡です。微分しても形が変わらない。つまり「今の大きさ」を微分した結果が「今の大きさ」になる。この自己相似的な美しさこそが、指数関数の本質です。

そして、この $e^x$ の逆、自然対数 $\ln x$ を微分すると $1/x$ になります。これは「増加の割合」が「元の大きさの逆数」に比例するという、極めて純粋な数学的構造を露わにします。

形式的な計算を捨て去り、この「成長と変化の調和」を感じ取ってください。これが微積分の真の醍醐味です！

\section{1. 「真の複利法則（True Compound Interest）」：成長し続ける元金}

トンプソンは第9章において、微積分学における最も神秘的で美しい定数――ネイピア数 $e = 2.71828\dots$ の本質を解き明かすために、経済学や銀行の「複利（Compound Interest）」のたとえ話を用います。

銀行に元金 $1$ 万円を預けたとしましょう。もし年利 $100\%$ という破格の銀行があったとします。

1年後に利息がまとめて支払われる「単利」であれば、1年後には元金 $1$ 万円と利息 $1$ 万円を合わせて、

$$1 + 1 = 2 \text{ 万円}$$

になります。

しかし、もし銀行が「半年ごと（年2回）」に利息を計算し、途中で生じた利息にもさらに利息をつけてくれる（半年複利）としたらどうなるでしょうか。半年ごとに $50\%$（0.5）の利息がつくので、1年後には、

$$\left(1 + \frac{1}{2}\right)^2 = (1.5)^2 = 2.25 \text{ 万円}$$

に増えます。

では、もし「毎月（年12回）」利息をつけてくれるなら？

$$\left(1 + \frac{1}{12}\right)^{12} \approx 2.613 \text{ 万円}$$

さらに「毎日（年365回）」利息を再投資してくれるなら？

$$\left(1 + \frac{1}{365}\right)^{365} \approx 2.7145 \text{ 万円}$$

利息の計算回数を増やせば増やすほど、金額は増えていきます。では、もし「毎秒、毎瞬、絶え間なく連続的に」利息が元金に組み込まれて自己増殖していくとしたら、1年後には資産は無限大に爆発してしまうでしょうか？

答えは「ノー」です。金額はある一定の限界値に向かってどこまでも近づいていきます。その極限値こそが、

$$\lim_{n \to \infty} \left(1 + \frac{1}{n}\right)^n = e = 2.718281828459\dots$$

という不思議な超越数、ネイピア数（自然対数の底）なのです。

\section{2. ネイピア数 $e = 2.71828\dots$ の発見：分割を無限にしたときの極限}

トンプソンは、この連続複利の極限プロセスを「真の複利法則（On True Compound Interest）」と呼び、自然界のあらゆる生物の増殖や物理現象の根底にある基本メカニズムであると見抜きました。

細菌の繁殖を考えてみてください。バクテリアは1年や1日の終わりに号令をかけて一斉に分裂するのではありません。すでに存在する個体の一匹一匹が、24時間絶え間なくリアルタイムに増殖し続けています。

放射性物質の崩壊も同じです。原子の集団は、今存在する原子の数に比例して、毎秒連続的に崩壊していきます。

コンデンサの放電、熱いコーヒーが冷めていく温度変化、大気圧の高度による減少――これらはすべて、「現在の量そのものに比例したペースで変化する」という共通の法則に従っています。

数式で言えば、変化率 $\frac{dy}{dt}$ が、現在の量 $y$ そのものに等しいという関係です：

$$\frac{dy}{dt} = y$$

「成長のスピードが、自分自身の大きさに等しい」――この自然界の最も根源的な要求を満たす唯一無二の数学的実体こそが、指数関数 $y = e^t$ なのです。

\section{3. 微分しても自分自身のまま：唯一無二の関数 $y = e^x$ の驚異}

指数関数 $y = e^x$ を微分してみましょう。微積分の教科書でこの性質に初めて出会ったとき、全ての学生が息を呑むような感動を覚えます。

$$\frac{d}{dx}\left(e^x\right) = e^x$$

なんと、微分しても形がまったく変わらず、自分自身のままなのです！

なぜこのような奇跡が起きるのでしょうか。トンプソン流の素朴な増分計算で確かめてみましょう。

$x$ が微小量 $dx$ だけ増えたとき、$y = e^x$ の増分は、

$$e^{x + dx} - e^x = e^x \cdot e^{dx} - e^x = e^x (e^{dx} - 1)$$

となります。

ここで、ネイピア数 $e$ の定義に基づき、$dx$ が極めて小さいとき、$e^{dx} \approx 1 + dx$ と展開されます（1次微小量の近似）。

したがって、

$$e^{x + dx} - e^x \approx e^x (1 + dx - 1) = e^x \, dx$$

両辺を $dx$ で割れば、

$$\frac{d}{dx}(e^x) = e^x$$

が見事に導かれます。

一般の底 $a$ を持つ指数関数 $y = a^x$（例えば $2^x$ や $10^x$）を微分すると、余計な定数係数 $\ln a$ が前に飛び出してきます（$\frac{d}{dx}(a^x) = a^x \ln a$）。

しかし、底が $e$ のときだけは、$\ln e = 1$ となるため、余計な係数が一切現れず、完全に純粋な自己保存が行われます。

物理学者やエンジニアが、あらゆる場面で $10^x$ ではなく $e^x$ を好んで使う理由はここにあります。ネイピア数 $e$ とは、微積分という演算に対して「最も自然で、摩擦のない底」なのです。

\section{4. 自然対数 $y = \ln x$ の微分：なぜ $\frac{d}{dx}(\ln x) = \frac{1}{x}$ になるのか}

指数関数 $y = e^x$ の相棒であり、その逆関数であるのが「自然対数（Natural Logarithm）」$y = \ln x$（または $\log_e x$）です。

この自然対数の微分公式もまた、驚くほど美しく、かつ極めて重要な意味を持っています。

$$\frac{d}{dx}(\ln x) = \frac{1}{x}$$

この導出は、第6章で学んだトンプソンの「連鎖律」と逆関数の関係を使えば、驚くほど簡単です。

$y = \ln x$ とは、定義により

$$x = e^y$$

と同値です。

両辺を $y$ で微分してみましょう。右辺の $e^y$ は微分しても $e^y$ のままですから、

$$\frac{dx}{dy} = e^y$$

となります。ここで $e^y = x$ でしたから、$\frac{dx}{dy} = x$ です。

私たちが求めたいのは $\frac{dy}{dx}$ ですから、分数の逆数をとるようにひっくり返せば、

$$\frac{dy}{dx} = \frac{1}{\frac{dx}{dy}} = \frac{1}{x}$$

が鮮やかに得られます！

この結果の歴史的意義は計り知れません。

第4章で学んだベキ関数の積分を思い出してください。$\int x^n dx = \frac{1}{n+1}x^{n+1}$ という公式において、$n = -1$ のときだけは分母が $n + 1 = 0$ となってしまい、計算が破綻していました。

「反比例の曲線 $y = \frac{1}{x}$ の面積を求める関数は何なのか？」という数学界の長年の謎に対して、まさにこの自然対数 $y = \ln x$ が、唯一無二の救世主として名乗りを上げたのです。

指数と対数は、自然界の成長と減衰を語る上で欠かすことのできない、微積分の両輪なのです。

\begin{examplebox}{高校生のための重要例題：第9章 —— 自己増殖する指数関数と自然対数の微分}
\textbf{【問題】}\\
関数 $y = x^2 e^x$ および $y = \ln(x^2 + 1)$ の導関数を求めよ。
\\[2mm]
\textbf{【詳細な解答・解説】}\\
(1) 積の微分公式 $(uv)' = u'v + uv'$ を適用する。$u = x^2, v = e^x$ とおくと、$u' = 2x, v' = e^x$（$e^x$ の微分は自分自身のまま）。\\$\frac{dy}{dx} = 2x e^x + x^2 e^x = (x^2 + 2x)e^x = x(x+2)e^x$。\\(2) 自然対数の合成関数である。連鎖律を用いる。$u = x^2 + 1$ とおくと $y = \ln(u)$。\\$\frac{dy}{du} = \frac{1}{u}, \quad \frac{du}{dx} = 2x$\\したがって $\frac{dy}{dx} = \frac{1}{u} \cdot 2x = \frac{2x}{x^2 + 1}$。
\end{examplebox}


\section*{第9章の確認演習}
\addcontentsline{toc}{section}{第9章の確認演習}
\begin{enumerate}[label=\textbf{問\arabic*.}]
\item 連続複利の極限 $\lim_{n\to\infty}\left(1 + \frac{1}{n}\right)^n$ からネイピア数 $e$ が定義される過程を説明せよ。
\item 関数 $y = e^x$ の微分がなぜ自分自身になるのか、その直観的証明を述べよ。
\item 自然対数 $y = \ln x$ の導関数が $\frac{1}{x}$ となることを、逆関数の微分法から示せ。
\end{enumerate}


\chapter{振動と三角関数の微分：波と円運動を描く $\sin x$ と $\cos x$ の導関数}
\chaptermark{振動と三角関数の微分}

\begin{flushright}
\small\textit{How to Deal with Sines and Cosines}
\end{flushright}

\begin{figure}[htbp]
\centering
\begin{tikzpicture}[x=1cm, y=1cm, >=stealth]
    % Axes
    \draw[->, thick] (-3.4, 0) -- (3.6, 0) node[right] {\footnotesize $\cos\theta$};
    \draw[->, thick] (0, -3.2) -- (0, 3.6) node[above] {\footnotesize $\sin\theta$};
    
    % Circle radius 2.6
    \draw[thick, darkslate] (0, 0) circle (2.6);
    
    % Angle theta = 45 deg
    % P = (2.6 cos 45, 2.6 sin 45) = (1.838, 1.838)
    \coordinate (P) at (1.838, 1.838);
    \draw[very thick, primaryblue] (0, 0) -- (P);
    \node[below right=2pt, primaryblue] at (0.9, 0.8) {\footnotesize 半径 $r=1$};
    \filldraw[primaryblue] (P) circle (2pt) node[right=3pt] {$P(\cos\theta, \sin\theta)$};
    
    % Angle arc
    \draw[->, thick, warmamber] (0.7, 0) arc (0:45:0.7);
    \node[warmamber] at (0.95, 0.4) {\footnotesize $\theta$};
    
    % Velocity vector (tangent, length 1.6)
    % VEnd = (0.707, 2.969)
    \coordinate (VEnd) at (0.707, 2.969);
    \draw[->, very thick, crimson] (P) -- (VEnd);
    % Velocity vector label inside triangle, below-left of tangent arrow
    \node[below left=2pt, crimson] at (1.25, 2.45) {\bfseries\footnotesize 速度ベクトル $\vec{v}$};
    
    % Components of velocity vector:
    % 1. Vertical component: up from P: (1.838, 1.838) -> (1.838, 2.969)
    \draw[thick, dashed, crimson] (P) -- (1.838, 2.969);
    \node[right=2pt, crimson] at (1.838, 2.4) {\footnotesize 垂直：$\cos\theta$};
    % 2. Horizontal component: left from (1.838, 2.969) -> (0.707, 2.969)
    \draw[thick, dashed, crimson] (1.838, 2.969) -- (VEnd);
    \node[above=2pt, crimson] at (1.27, 2.98) {\footnotesize 水平：$-\sin\theta$};
    
    % Bottom Summary Box
    \node[draw=primaryblue, fill=softblue, rounded corners=4pt, text width=11.0cm, align=center, inner sep=5pt] at (0, -3.7) {
        \textbf{三角関数の導関数}：$\frac{d}{d\theta}(\sin\theta) = \cos\theta$、\quad $\frac{d}{d\theta}(\cos\theta) = -\sin\theta$
    };
\end{tikzpicture}
\caption{単位円上の等速円運動における位置と速度ベクトルの直交性・三角関数の導関数}
\label{fig:thompson_trig_circ}
\end{figure}

教育情報分析研究会　数学教育・物理解析リサーチユニット

\section{【Ollama専任講師特別講義】トンプソン流直観的微積分の扉}

皆さん、ようこそ！今日は数学の「真の美しさ」に触れる準備はいいですか？

私たちはこれまで、公式を暗記するだけの「作業」をしてきたかもしれません。しかし、トンプソンは『Calculus Made Easy』において、微積分を「変化の幾何学」として捉えることを説いています。今日はその核心、三角関数の微分を解剖しましょう。

まず、なぜ弧度法（ラジアン）なのか。これは単なる習慣ではありません。度数法は「円を360分割する」という人間の恣意的な約束ですが、ラジアンは「半径と弧の長さが等しくなる角度」という、円の幾何学的性質そのものから導かれる自然な単位です。これを使うことで、微分の極限において係数が「1」という最も純粋な形に収束するのです。

では、$\sin x$ の微分が $\cos x$ になる理由を、単位円から見てみましょう。単位円上の点 $(x, y)$ において、$y = \sin \theta$ です。この点における接線の傾き（微分）を求めると、幾何学的に「その地点における水平方向の速度に対する垂直方向の速度の比」となります。単位円の円周方向の動きを解析すると、$\sin$ の変化率は、その瞬間の水平位置 $\cos$ そのものに一致することが導かれます。これが「$\sin x$ を微分すると $\cos x$」の幾何学的本質です。

さらに、これを2回微分すると $\sin x$ にマイナスがつきます。これは物理における「単振動」の正体です。加速（2階微分）が位置に比例し、かつ向きが逆（マイナス）であるという関係は、復元力が常に中心へ向かうことを意味します。

公式の裏にある「幾何学的な必然性」を感じ取ってください。数学は、この美しい一貫性の上に成り立っているのです！

\section{1. なぜ角度を度（度数法）ではなくラジアン（弧度法）で測るのか}

トンプソンは第10章において、周期的な運動（振り子の揺れ、音波、交流電流、光の波）を記述する必須の道具である「三角関数（正弦 $\sin$ と余弦 $\cos$）」の微分を取り上げます。そして冒頭で、高校生が必ずぶつかる疑問――「なぜ微積分では、日常使い慣れた度数法（$360^\circ$）を捨てて、ラジアン（弧度法）を使わなければならないのか？」という謎を明快に解き明かします。

円を360等分する度数法は、古代バビロニアの天文学に由来する人為的な約束事に過ぎません。もし度数法のまま $\sin x$ を微分しようとすると、計算の途中で $\frac{\pi}{180}$ という厄介な換算係数が毎回飛び出してきて、数式が極めて見苦しくなってしまいます。

これに対し、「ラジアン（弧度法）」は、半径 $r$ の円において「弧の長さがちょうど半径 $r$ と同じになる角度」を $1$ ラジアンと定義する、純粋に幾何学的で自然な単位系です。

円周全体の長さは $2\pi r$ ですから、1周360度は $2\pi$ ラジアン（約6.28ラジアン）になります。

この弧度法を採用することで、微積分の世界において最も美しく重要な極限の等式が誕生します：

$$\lim_{\theta \to 0} \frac{\sin\theta}{\theta} = 1$$

微小な角度 $\theta$ においては、「角度そのものの値（ラジアン）」と「その正弦 $\sin\theta$」が、1次微小量として完全に一致する（$\sin\theta \approx \theta$）のです。この奇跡的な一致こそが、三角関数の微積分を摩擦なく滑らかに回転させるための絶対条件なのです。

\section{2. 直観的幾何学証明：単位円上の回転と微小直角三角形}

トンプソンは、$\sin x$ の微分公式を、抽象的な加法定理の式変形だけに頼るのではなく、単位円（半径 $1$ の円）を用いた鮮やかな幾何学的直観によって示します。

半径 $1$ の円の円周上を反時計回りに動く点 $P$ を考えましょう。$x$ 軸の正の向きからの回転角を $\theta$ とすると、点 $P$ の座標は、

$$x = \cos\theta, \quad y = \sin\theta$$

となります。

いま、角度がほんの少しだけ $d\theta$ だけ回転したとします。点 $P$ は円周に沿って新しい点 $Q$ へ移動します。

円の半径は $1$ ですから、弧 $PQ$ の長さは、弧度法の定義そのものにより、

$$\text{弧 } PQ = 1 \cdot d\theta = d\theta$$

となります。

$d\theta$ は微小量ですから、この極小の弧 $PQ$ は、円の接線方向を向いた「長さ $d\theta$ の真っ直ぐな線分」とみなせます。

点 $P$ から点 $Q$ への微小な移動を、水平方向の成分 $dx$ と垂直方向の成分 $dy$ に分解して、直角三角形を作りましょう。

円の接線は半径と直交しますから、初等幾何学の角度の追跡により、この微小直角三角形の斜辺が $d\theta$、高さが $dy$、底辺が $-dx$ となり、その内角が元の角度 $\theta$ とぴったり一致することがわかります。

したがって、直角三角形の三角比から直ちに次の関係が得られます：

$$\frac{dy}{d\theta} = \cos\theta$$

$y = \sin\theta$ でしたから、これはまさに「$\sin\theta$ の導関数は $\cos\theta$ である」という結論を幾何学的に証明しています！

\section{3. 基本極限 $\lim_{\theta \to 0} \frac{\sin\theta}{\theta} = 1$ と $(\sin x)' = \cos x$}

この幾何学的直観を、三角関数の加法定理を用いた代数的な微小量計算で確かめてみましょう。

$\theta$ を $x$ と書き換えて、関数 $y = \sin x$ の微小増分を計算します。

$x$ を $x + dx$ に増やすと、

$$\sin(x + dx) = \sin x \cos(dx) + \cos x \sin(dx)$$

となります。

ここで、$dx$ は極小量ですから、

\begin{itemize}\item $\cos(dx) \approx 1$（2次微小量 $(dx)^2$ の誤差を無視）\end{itemize}

\begin{itemize}\item $\sin(dx) \approx dx$（基本極限による1次微小量）\end{itemize}

と置き換えることができます。

これを代入すると、

$$\sin(x + dx) \approx \sin x \cdot 1 + \cos x \cdot dx = \sin x + \cos x \, dx$$

元の $\sin x$ を差し引くと、正味の微小増分は

$$dy = \cos x \, dx$$

となります。両辺を $dx$ で割れば、

$$\frac{d}{dx}(\sin x) = \cos x$$

が得られます。

なんと淀みのない論理でしょう！

正弦関数 $\sin x$ の波を微分すると、位相が $\frac{\pi}{2}$（90度）だけ進んだ余弦関数 $\cos x$ の波にそっくりそのまま変身するのです。

サインのグラフの頂点（$x = \pi/2$）では接線が水平（傾きゼロ）になりますが、確かに $\cos(\pi/2) = 0$ となって整合しています。

\section{4. 余弦の微分 $(\cos x)' = -\sin x$ と単振動・波動への展開}

次に、余弦関数 $y = \cos x$ の微分を調べましょう。先ほどと同様に微小増分を計算すると、

$$\cos(x + dx) = \cos x \cos(dx) - \sin x \sin(dx) \approx \cos x \cdot 1 - \sin x \, dx$$

となります。増分は $dy = -\sin x \, dx$ ですから、導関数は

$$\frac{d}{dx}(\cos x) = -\sin x$$

となります。ここにマイナス符号が現れることに注目してください！

$x = 0$ から出発するとき、$\cos x$ は最大値 $1$ から減少していくため、傾きは当然マイナスでなければなりません。

ここで、$\sin x$ を2回微分してみましょう（第7章の高階微分）：

\begin{enumerate}\item 1回微分：$\frac{d}{dx}(\sin x) = \cos x$（速度）\end{enumerate}

\begin{enumerate}\item 2回微分：$\frac{d^2}{dx^2}(\sin x) = \frac{d}{dx}(\cos x) = -\sin x$（加速度）\end{enumerate}

なんと、2回微分すると、「元の関数にマイナスをつけたもの」に戻ってきました！

$$\frac{d^2y}{dx^2} = -y$$

この数式こそ、物理学における最も重要な方程式のひとつ、単振動の方程式（Harmonic Oscillator）です。

バネに繋がれたおもりの振動、振り子の往復運動、バイオリンの弦の響き、電波のアンテナから放射される電磁波――これら自然界のあらゆる「振動と波」は、この第2次導関数が元の位置の逆方向に向く（復元力が働く）という微分方程式に支配されています。トンプソンが説いた三角関数の微積分は、宇宙を満たすあらゆる波のシンフォニーを解明するための鍵なのです。

\begin{examplebox}{高校生のための重要例題：第10章 —— 三角関数の微分と合成関数}
\textbf{【問題】}\\
次の各関数を微分せよ。\\(1) $y = \sin(3x)$ \quad (2) $y = \cos^2(x)$ \quad (3) $y = x \sin(x)$
\\[2mm]
\textbf{【詳細な解答・解説】}\\
(1) 連鎖律：$u = 3x$ とおくと $y = \sin(u)$。$\frac{dy}{dx} = \cos(u) \cdot 3 = 3\cos(3x)$。\\(2) 連鎖律：$u = \cos(x)$ とおくと $y = u^2$。$\frac{dy}{dx} = 2u \cdot (-\sin x) = -2\sin(x)\cos(x) = -\sin(2x)$。\\(3) 積の微分：$u = x, v = \sin(x)$。$\frac{dy}{dx} = 1 \cdot \sin(x) + x \cdot \cos(x) = \sin(x) + x\cos(x)$。
\end{examplebox}


\section*{第10章の確認演習}
\addcontentsline{toc}{section}{第10章の確認演習}
\begin{enumerate}[label=\textbf{問\arabic*.}]
\item 度数法ではなく弧度法（ラジアン）を用いたときにのみ $\lim_{\theta\to 0}\frac{\sin\theta}{\theta} = 1$ が成立する理由を説明せよ。
\item 等速円運動の射影から、$\frac{d}{dx}(\sin x) = \cos x$ と $\frac{d}{dx}(\cos x) = -\sin x$ を導け。
\item 単振動の方程式 $\frac{d^2y}{dx^2} = -\omega^2 y$ の解が三角関数となる理由を考察せよ。
\end{enumerate}


\chapter{微分の逆演算としての積分：不定積分の基礎と原始関数の探求}
\chaptermark{微分の逆演算としての積分}

\begin{flushright}
\small\textit{Integration: Reversing the Process and Finding Solutions}
\end{flushright}

\begin{figure}[htbp]
\centering
\begin{tikzpicture}[x=1cm, y=1cm, >=stealth]
    % Outer Frame
    \draw[thick, fill=softrose!30, draw=crimson, rounded corners=6pt] (-7.6, -4.2) rectangle (7.6, 4.0);
    \node[draw=crimson, fill=white, rounded corners=4pt, inner sep=5pt] at (0, 3.2) {\bfseries 微分と積分の双対性：過程の逆行（Reversing the Process）};
    
    % Left Box: Primitive Functions
    \node[draw=primaryblue, fill=white, rounded corners=4pt, text width=4.0cm, align=center, inner sep=6pt] at (-4.8, 0.9) {
        \textbf{原始関数族}\\[3pt]
        \Large $F(x) + C$\\[3pt]
        \footnotesize 積分定数 $C$ の任意性\\[2pt]
        \footnotesize 例：$\frac{1}{3}x^3 + C$\\[3pt]
        \scriptsize 「元の滑らかな全体」
    };
    
    % Right Box: Derivatives
    \node[draw=accentgreen, fill=white, rounded corners=4pt, text width=4.0cm, align=center, inner sep=6pt] at (4.8, 0.9) {
        \textbf{導関数 $f(x)$}\\[3pt]
        \footnotesize 瞬間の傾き・変化率\\[3pt]
        \footnotesize 例：$x^2$\\[4pt]
        \scriptsize 「微小なかけらの勢い」
    };
    
    % Center Arrows with clear labels (gap: -2.8 to 2.8 = 5.6cm)
    \draw[->, very thick, primaryblue] (-2.4, 1.4) -- (2.4, 1.4);
    \node[above=3pt] at (0, 1.4) {\bfseries\small 微分演算 $\frac{d}{dx}$（定数 $C$ の消失）};
    
    \draw[<-, very thick, crimson] (-2.4, 0.4) -- (2.4, 0.4);
    \node[below=3pt] at (0, 0.4) {\bfseries\small 不定積分 $\int \dots dx$（定数 $C$ の創出）};
    
    % Bottom Formula Box
    \node[draw=darkslate, fill=white, rounded corners=4pt, text width=13.6cm, align=center, inner sep=6pt] at (0, -2.6) {
        \textbf{基本積分公式}：\Large $\int x^n dx = \frac{1}{n+1}x^{n+1} + C \quad (n \neq -1)$、\quad $\int \frac{1}{x}dx = \ln|x| + C \quad (n = -1)$
    };
\end{tikzpicture}
\caption{微分の逆演算としての不定積分と積分定数 $C$ の幾何学的役割}
\label{fig:thompson_integration_reverse}
\end{figure}

教育情報分析研究会　数学教育・物理解析リサーチユニット

\section{【Ollama専任講師特別講義】トンプソン流直観的微積分の扉}

皆さん、ようこそ。今日は、トンプソンが提唱した「微積分の真の精神」の核心、第11章「不定積分」へと踏み込みます。

いいですか、微積分とは「変化のルール」を操る学問です。微分が「ある瞬間の変化率（導関数）」を求める操作なら、積分はその「逆操作」です。トンプソンはこれを、時計の針を巻き戻すような、極めて論理的な「推理の旅」として描き出しました。

「この変化のルール（導関数）から、元の関数（原始関数）を導き出せ」という問いに対し、我々は数学的な足跡を辿り、元の形を復元するのです。

ここで登場する長い「S」の記号（$\int$）。これは単なる記号ではありません。それは「積み重ね（Sum）」の略称であり、微小な変化を無限に積み上げて元の全体像を再構成するという、数学的意志の象徴です。

そして、最も重要な「消えた手がかり」――積分定数 $C$。微分という操作は、関数の「傾き」を抽出する過程で、元の関数の「高さ（初期位置）」の情報を削り取ってしまいます。積分で元の関数を復元しようとしても、その「削り取られた初期位置」までは特定できない。だからこそ、私たちは $C$ という「不確定な空白」を認めるのです。

この $C$ は、数学的な誠実さの証です。元の関数が一つに決まらないという事実を正しく受け入れること。それこそが、トンプソンが教えようとした微積分の本質なのです。

\section{1. 時計の針を巻き戻す：「微分する前の元の関数」を探す旅}

これまでの10章にわたって、私たちは「ある関数が与えられたとき、その微小なかけらを切り出し、変化率（導関数）を求める」という「微分（Differentiation）」の技術を徹底的に学んできました。

第11章において、トンプソンは読者を微積分のもうひとつの巨大な柱――「積分（Integration）」の世界へと導きます。

トンプソンは積分の導入として、意表を突くような平易な問いを投げかけます。

「もし誰かがあなたに、ある関数の導関数を見せて、『私が微分する前の、元の関数は何だったか当ててごらん』と言ったら、あなたはどう答えるだろうか？」

例えば、「微分したら $2x$ になった」と言われれば、私たちは第4章の知識から即座に「元の関数は $x^2$ だ！」と答えることができます。

「微分したら $3x^2$ になった」と言われれば、「元の関数は $x^3$ だ！」と分かります。

このように、「微分のプロセスを逆向きに巻き戻し、導関数から元の関数（原始関数、Primitive function）を推理・復元する操作」のことを、数学では不定積分（Indefinite Integration）と呼びます。

微分が「暗号化」だとすれば、積分は「復号（暗号解読）」です。微分が「細かく砕くこと」だとすれば、積分は「組み立てて元に戻すこと」です。微分の公式をマスターした皆さんにとって、不定積分の基本的なルールは、単に微分の公式を右から左へ逆再生することに他なりません。

\section{2. 記号 $\int$（長いS）の登場：微小量 $f(x)dx$ を集めて復元する}

第1章で予告された、あの細長い記号「$\int$」が、ここで本格的に主役に躍り出ます。

トンプソンは、積分の記号表現の構造を丁寧に解剖します：

$$\int f(x)\,dx$$

この式を観察してください。単に $\int$ という記号だけがあるのではなく、後ろに必ず「$dx$」がセットで寄り添っています。なぜでしょうか。

もし $F(x)$ を微分した結果が $f(x)$ だったとしましょう：

$$\frac{dF}{dx} = f(x)$$

両辺に微小量 $dx$ を掛け算して払ってみると、

$$dF = f(x)\,dx$$

となります。

この式の右辺 $f(x)\,dx$ は、「元の関数 $F$ の、ほんの少しの微小なかけら $dF$」を表しています。

そして記号 $\int$ は、ラテン語の *summa*（総和、寄せ集め）を意味する長いSでした。

したがって、両辺に $\int$ を作用させるということは、

$$\int dF = \int f(x)\,dx$$

すなわち、「微小なかけら $dF$ を、全部かき集めて足し合わせる」という操作を意味します。

かけら $dF$ を全部集めれば、当然ながら元の全体量 $F(x)$ が復元されます！

これが、不定積分の記号 $\int f(x)\,dx$ の直観的構造です。積分とは、微小なかけら $f(x)dx$ を足し合わせて元の全体を取り戻す作業なのです。

\section{3. 積分定数 $C$ の謎：なぜ微分すると定数は消え、積分すると現れるのか}

しかし、時計の針を巻き戻すこの推理の旅には、ひとつだけ重大な落とし穴があります。それこそが、高校生を大いに悩ませる「積分定数 $C$（Constant of Integration）」の存在です。

次の3つの関数を微分してみてください：

\begin{enumerate}\item $y_1 = x^2$\end{enumerate}

\begin{enumerate}\item $y_2 = x^2 + 5$\end{enumerate}

\begin{enumerate}\item $y_3 = x^2 - 100$\end{enumerate}

それぞれの導関数を計算すると、

$$\frac{dy_1}{dx} = 2x, \quad \frac{dy_2}{dx} = 2x + 0 = 2x, \quad \frac{dy_3}{dx} = 2x - 0 = 2x$$

なんと、3つの関数はどれも全く同じ導関数 $2x$ になってしまいます！

なぜなら、定数項（固定された数値）は $x$ がいくら変化しても全く動かない（変化率がゼロ）ため、微分という操作を通過すると完全に消滅して蒸発してしまうからです。

ということは、逆に「微分した結果が $2x$ だった」という情報だけから元の関数を推理するとき、元の関数が $x^2$ だったのか、$x^2 + 5$ だったのか、$x^2 - 100$ だったのか、決定することができません。

定数の部分が何であったかは、微分によって完全に闇の中に消え去ってしまったからです。

そこで、数学者はこの「消えてしまったかもしれない任意の定数」を表すために、末尾に「$+ C$」という記号を添える約束をしました：

$$\int 2x\,dx = x^2 + C$$

この $C$ を積分定数と呼びます。

トンプソンはこの積分定数を、犯罪捜査における「消えた手がかり」に例えています。現場に残された足跡（導関数 $2x$）だけでは犯人の身長（グラフの形状）は特定できても、犯人が何階の部屋に住んでいるか（上下のシフト位置）までは分かりません。部屋の階数を特定するためには、「$x = 0$ のとき $y = 3$ である」といった追加の目撃証言（初期条件・境界条件）が必要になるのです。

\section{4. 逆ベキの公式 $\int x^n dx = \frac{1}{n+1}x^{n+1} + C$ と例外 $n = -1$（対数の登場）}

微分のベキ関数の公式 $\frac{d}{dx}(x^m) = m x^{m-1}$ を逆再生することで、不定積分の最も基本的で強力な公式が直ちに導かれます。

微分すると指数が $1$ つ減るのであれば、積分するときは指数を $1$ つ増やさなければなりません。

$x^{n+1}$ を微分すると $(n+1)x^n$ になりますから、余計な係数 $(n+1)$ を相殺するために、最初から $\frac{1}{n+1}$ で割っておけばよいのです。

かくして、逆ベキの積分公式が完成します：

$$\int x^n \, dx = \frac{1}{n+1} x^{n+1} + C \quad (n \neq -1)$$

例えば：

\begin{itemize}\item $\int x^2 dx = \frac{1}{3}x^3 + C$\end{itemize}

\begin{itemize}\item $\int x^5 dx = \frac{1}{6}x^6 + C$\end{itemize}

\begin{itemize}\item $\int 1 \, dx = \int x^0 dx = x + C$\end{itemize}

\begin{itemize}\item $\int \frac{1}{x^2} dx = \int x^{-2} dx = \frac{1}{-1}x^{-1} + C = -\frac{1}{x} + C$\end{itemize}

\begin{itemize}\item $\int \sqrt{x} dx = \int x^{1/2} dx = \frac{1}{3/2}x^{3/2} + C = \frac{2}{3}x\sqrt{x} + C$\end{itemize}

ここで、公式の横に添えられた条件「$n \neq -1$」に刮目してください。

もし $n = -1$（すなわち $\int \frac{1}{x} dx$）でこの公式を使おうとすると、分母が $n + 1 = -1 + 1 = 0$ となってしまい、ゼロ除算の禁忌に触れて破綻してしまいます。

では、$\frac{1}{x}$ の不定積分は何になるのでしょうか。

第9章を思い出してください！ 私たちは「微分したら $\frac{1}{x}$ になる関数」をすでに知っています。それこそが、自然対数 $\ln x$ でした！

$$\int \frac{1}{x}\,dx = \ln|x| + C$$

欠けていた最後のパズルのピースが、ここで美しく嵌まりました。多項式から分数関数、無理関数に至るまで、積分の道具立ては完璧に整ったのです。

\begin{examplebox}{高校生のための重要例題：第11章 —— 不定積分の計算と原始関数の決定}
\textbf{【問題】}\\
次の不定積分を求めよ。\\(1) $\int (6x^2 - 4x + 5)\,dx$ \quad (2) $\int \left(\frac{1}{x} + \frac{2}{\sqrt{x}}\right)dx$ \quad (3) $\int 3e^x dx$
\\[2mm]
\textbf{【詳細な解答・解説】}\\
(1) ベキ関数の基本積分公式 $\int x^n dx = \frac{1}{n+1}x^{n+1} + C$ より、\\$\int (6x^2 - 4x + 5)\,dx = 6\left(\frac{x^3}{3}\right) - 4\left(\frac{x^2}{2}\right) + 5x + C = 2x^3 - 2x^2 + 5x + C$。\\(2) 指数表記で整理すると $\int (x^{-1} + 2x^{-1/2})dx$。\\$x^{-1}$ の積分は自然対数 $\ln|x|$ となり、$x^{-1/2}$ の積分は $\frac{x^{1/2}}{1/2} = 2\sqrt{x}$ であるから、\\$\int \left(\frac{1}{x} + \frac{2}{\sqrt{x}}\right)dx = \ln|x| + 4\sqrt{x} + C$。\\(3) 指数関数の積分は自分自身のままであるから、$\int 3e^x dx = 3e^x + C$。（すべて $C$ は積分定数）。
\end{examplebox}


\section*{第11章の確認演習}
\addcontentsline{toc}{section}{第11章の確認演習}
\begin{enumerate}[label=\textbf{問\arabic*.}]
\item 「積分とは過程の逆行である」というトンプソンの思想と、積分定数 $C$ の必然性を論ぜよ。
\item なぜ $\int x^{-1} dx$ においてベキ関数の積分公式が使えず、自然対数が出現するのか説明せよ。
\item 置換積分法と部分積分法の公式を、微分の合成関数・積の微分公式から導出せよ。
\end{enumerate}


\chapter{面積の求積と微積分学の基本定理：区分求積法から現代科学・AIへの架け橋}
\chaptermark{面積の求積と基本定理}

\begin{flushright}
\small\textit{On Finding Areas by Integrating and the Fundamental Theorem}
\end{flushright}

\begin{figure}[htbp]
\centering
\begin{tikzpicture}[x=1cm, y=1cm, >=stealth]
    % Top Header Box
    \node[draw=warmamber, fill=softamber, rounded corners=4pt, inner sep=5pt] at (4.0, 4.6) {
        \bfseries 微積分学の基本定理：$\frac{dS}{dx} = f(x) \iff \int_a^b f(x)dx = F(b) - F(a)$
    };
    
    % Axes
    \draw[->, thick] (0, 0) -- (8.2, 0) node[right] {$x$};
    \draw[->, thick] (0, 0) -- (0, 3.8) node[left] {$y$};
    
    % Curve y = 0.04 x^2 + 0.12 x + 0.8
    \draw[very thick, primaryblue, domain=0.8:7.2, samples=60] plot (\x, {0.04*\x*\x + 0.12*\x + 0.8});
    \node[above left, primaryblue] at (6.8, 3.6) {\bfseries\footnotesize 曲線 $y = f(x)$};
    
    % a=1.8, b=6.2
    \coordinate (A) at (1.8, 0);
    \coordinate (B) at (6.2, 0);
    
    % Shaded Area S(x) from a=1.8 to x=3.6
    \fill[softblue, opacity=0.8] (1.8, 0) -- plot[domain=1.8:3.6] (\x, {0.04*\x*\x + 0.12*\x + 0.8}) -- (3.6, 0) -- cycle;
    \node at (2.7, 0.7) {\small\textbf{面積 $S(x)$}};
    
    % Strip dx at [3.6, 4.1]
    \fill[softamber] (3.6, 0) -- plot[domain=3.6:4.1] (\x, {0.04*\x*\x + 0.12*\x + 0.8}) -- (4.1, 0) -- cycle;
    \draw[thick, warmamber] (3.6, 0) -- (3.6, {0.04*3.6*3.6 + 0.12*3.6 + 0.8}) 
                         -- (4.1, {0.04*4.1*4.1 + 0.12*4.1 + 0.8}) -- (4.1, 0);
    
    % Strip label with clear arrow pointing from upper left
    \draw[<-, thick, warmamber] (3.85, 1.4) -- (2.6, 2.5);
    \node[above, warmamber] at (2.6, 2.5) {\bfseries\footnotesize 微小短冊 $dS \approx f(x)dx$};
    
    % Boundaries
    \draw[dotted, thick] (1.8, {0.04*1.8*1.8 + 0.12*1.8 + 0.8}) -- (1.8, 0) node[below] {$a$};
    \draw[dotted, thick] (6.2, {0.04*6.2*6.2 + 0.12*6.2 + 0.8}) -- (6.2, 0) node[below] {$b$};
    
    % Sub-labels on x axis
    \node[below=2pt] at (3.6, 0) {\footnotesize $x$};
    \node[below=14pt] at (4.1, 0) {\footnotesize $x+dx$};
\end{tikzpicture}
\caption{短冊の微小面積 $dS = f(x)dx$ による区分求積と微積分学の基本定理}
\label{fig:thompson_area_fundamental}
\end{figure}

教育情報分析研究会　数学教育・物理解析リサーチユニット

\section{【Ollama専任講師特別講義】トンプソン流直観的微積分の扉}

「皆さん、今日、数学の歴史における最大の『革命』の瞬間を共有しましょう。

微積分学の基本定理は、単なる計算テクニックではありません。それは『面積』と『変化率』という、一見無関係な二つの概念を一本の線で結んだ、人類史上最も美しい知の統合です。

かつてアルキメデスは、円周の面積を求めるために『取り尽くし法』を用いました。細かな断片を積み上げ、極限まで細分化する。しかし、この『足し合わせ』の作業は、非常に骨の折れるプロセスでした。

ここでトンプソンが突きつける数学的本質があります。それは『面積の増加率が、その瞬間の高さ $f(x)$ に等しい』という大発見です。

想像してください。面積を『積み上げる』のではなく、面積が『増えていく速度』に着目するのです。面積が動く速度（微分）を計算すれば、それが積み重なった総量（積分）を導き出せる。この逆転の発想こそが、積分と微分の繋がりを決定づけたのです。

この原理は現代のAI、特に深層学習の『勾配降下法』に直結しています。AIが正解へ向かうための最適解を求めるプロセスは、まさにこの微積分学の基本定理を数億回繰り返していることに他なりません。

微積分は、静止した形を測る学問から、動的な変化を支配する学問へと進化したのです。この感動を、皆さんの血肉にしてください！」

\section{1. 湾曲した図形の面積を求める古代からの夢：エウドクソスとアルキメデスの取り尽くし法}

トンプソンは『Calculus Made Easy』の掉尾を飾る第12章において、微積分学が到達した人類最高峰の智恵の結晶――「微積分学の基本定理（Fundamental Theorem of Calculus）」を語り尽くします。それは、古代ギリシャ以来2000年以上にわたって人類を魅了し、悩ませ続けてきた「曲線の面積を求める」という夢が、いかにして微分と奇跡の合流を果たしたかという壮大な物語です。

正方形や長方形、三角形のような真っ直ぐな線で囲まれた図形の面積なら、誰でも簡単に計算できます。しかし、ひとたび境界線が滑らかに曲がりくねった曲線（放物線、円、三次関数など）になった途端、初等幾何学の公式は無力化します。

古代ギリシャの大数学者エウドクソスやアルキメデス（Archimedes, 前287--前212）は、この難問に対して「取り尽くし法（Method of Exhaustion）」という執念の技法を編み出しました。

曲線の内側と外側に、無数の細長い長方形や三角形をぎっしりと敷き詰め、その短冊の数を増やしていくことで、真の面積を両側から挟み撃ちにして追い詰めるという方法です。

アルキメデスはこの方法で放物線の面積や円周率を計算することに成功しましたが、それは図形ごとに極めて高度で複雑なパズルを解くような、途方もない超人的努力を要するものでした。

\section{2. リーマン和と短冊の極限：$\sum f(x_i)\Delta x \to \int_a^b f(x) dx$}

この古代の取り尽くし法を、近代数学の普遍的なアルゴリズムへと昇華させたのが「定積分（Definite Integral）」です。

関数 $y = f(x)$ のグラフと $x$ 軸、および2本の直線 $x = a, x = b$ で囲まれた領域の面積 $S$ を求める問題を考えましょう。

区間 $[a, b]$ を $n$ 個の微小な小区間に等分割します。各区間の幅を $\Delta x = \frac{b - a}{n}$ と置きます。

各区間において、高さ $f(x_i)$、幅 $\Delta x$ の細長い長方形（短冊）を作ります。

1本の短冊の面積は、底辺 $\times$ 高さ、すなわち

$$\Delta S_i = f(x_i) \cdot \Delta x$$

です。

これら $n$ 本の短冊の面積を端から端まで全部足し合わせると、

$$S_n = \sum_{i=1}^n f(x_i) \Delta x$$

となります。これを「リーマン和（Riemann Sum）」と呼びます。階段状のギザギザがあるため、真の曲線の面積とはわずかな隙間（誤差）があります。

しかし、ここで分割数 $n$ を無限大に増やし、短冊の幅 $\Delta x$ を限りなくゼロに近づける極限（$\Delta x \to 0$）をとるとどうなるでしょうか。

短冊の幅は微小量 $dx$ となり、ギザギザの隙間は完全に消失します。

そして、ギリシャ文字の和の記号 $\sum$ は、あのライプニッツの滑らかな連続和の記号 $\int$ へと姿を変えます：

$$\lim_{n \to \infty} \sum_{i=1}^n f(x_i) \Delta x = \int_a^b f(x)\,dx$$

これが「定積分」の幾何学的定義です。定積分とは、幅 $dx$、高さ $f(x)$ の無限に細い短冊を、区間 $[a, b]$ にわたって隙間なく貼り合わせた「面積そのもの」なのです。

\section{3. 微積分学の基本定理：面積の増加率は関数の高さに等しい}

定積分の幾何学的意味は分かりましたが、アルキメデスのように無数の短冊の足し算を毎回極限計算していたのでは、計算が大変すぎて実用になりません。

17世紀後半、アイザック・ニュートンとライプニッツは、独立に、人類の歴史を永遠に変える戦慄すべき大発見を成し遂げました。

「面積を求める足し算（積分）は、なんと接線の傾きを求める割り算（微分）の、完全な逆演算になっている！」

この奇跡の法則を「微積分学の基本定理」と呼びます。トンプソンは、この定理の直観的証明を次のように鮮やかに解説します。

$x = a$ から任意の場所 $x$ までの面積を、関数 $S(x)$ とおきます。

いま、$x$ をほんの少しだけ $dx$ だけ右へ進めたとしましょう。

その結果、新しく増えた面積 $dS$ は、どのような形をしているでしょうか。

幅が $dx$ で、高さがグラフのその地点の高さ $f(x)$ である、ごく細い長方形の短冊1本分です！

$$dS \approx f(x) \cdot dx$$

両辺を底辺の幅 $dx$ で割ってみてください：

$$\frac{dS}{dx} = f(x)$$

なんと！ 「面積関数 $S(x)$ の導関数（成長の勢い）は、元の曲線の高さ $f(x)$ そのものになる」 のです！

これは決定的な事実を告げています。

面積 $S(x)$ を知りたければ、難解な短冊の極限和を計算する必要など一切ないのです。

「微分したら $f(x)$ になる関数（原始関数 $F(x)$）」を探し出しさえすればよいのです！

原始関数を $F(x)$ とすると、$S(x) = F(x) + C$ と書けます。$S(a) = 0$（幅ゼロの面積はゼロ）より $C = -F(a)$ ですから、区間 $[a, b]$ の全面積は、

$$\int_a^b f(x)\,dx = F(b) - F(a)$$

という、わずか引き算1回で瞬時に求まってしまうのです！

この瞬間、幾何学（面積）と代数学（微分）は完全に一体化し、人類は曲線の面積、立体の体積、曲線の長さ、物理的な仕事量を、自由自在に計算する神の如き計算力を手に入れたのです。

\section{4. 現代社会を支える微積分：物理シミュレーションからAIディープラーニングへ}

シルバヌス・P・トンプソンが『Calculus Made Easy』を通じて私たちに手渡してくれたこの知的遺産は、現代社会のありとあらゆるテクノロジーの心臓部で脈打っています。

高校物理で学ぶ力学において、力 $F(x)$ を位置 $x$ で積分すれば「仕事（エネルギー） $W = \int F dx$」となり、速度 $v(t)$ を時間 $t$ で積分すれば「移動距離 $x = \int v dt$」となります。

電気工学においては、電流 $i(t)$ を積分すれば蓄積された電荷 $q = \int i dt$ となり、現代のスマートフォンやコンピュータの回路設計は、マクスウェル方程式の微積分によって動いています。

航空宇宙工学では、ロケットのツィオルコフスキーの公式が微積分によって導かれ、気象予測の巨大スーパーコンピュータは、地球規模の大気の流体力学方程式を毎秒何兆回も数値積分しています。

そして21世紀、最も熱狂的な注目を集める人工知能（AI）と深層学習（ディープラーニング）の根底にあるアルゴリズム「勾配降下法（Gradient Descent）」や「誤差逆伝播法（Backpropagation）」は、まさに第3章の導関数 $\frac{dy}{dx}$ と第6章の連鎖律（チェイン・ルール）そのものです。

何億個、何兆個ものパラメータを持つ巨大なニューラルネットワークが、画像や言語を正しく認識できるように学習するプロセスとは、損失関数（誤差）の勾配（多変数の導関数）を計算し、極小値（第8章の最適化）へ向かってパラメータを少しずつ更新していくことに他なりません。

「What one fool can do, another can.」

一人の人間が理解できる美しい真理は、他のどんな人間にも必ず直観的に理解できる――トンプソンのこの確信は、100年以上の時を超えて、現代の私たちに力強く語りかけています。

微積分とは、一部の専門家のための難解な記号の牢獄ではありません。世界の変化を捉え、自然の調和を感じ取り、自らの手で未来のテクノロジーを創造するための、人類に等しく開かれた「自由の翼」なのです。

\begin{examplebox}{高校生のための重要例題：第12章 —— 定積分による放物線下の面積求積}
\textbf{【問題】}\\
放物線 $y = 3x^2$ と $x$ 軸、および2直線 $x = 1, x = 3$ で囲まれた図形の面積 $S$ を、定積分の定義（微積分学の基本定理）に基づき求めよ。
\\[2mm]
\textbf{【詳細な解答・解説】}\\
求める面積 $S$ は、微小短冊 $y \cdot dx = 3x^2 dx$ を $x = 1$ から $x = 3$ まで連続的に足し合わせた定積分として表される：\\$S = \int_1^3 3x^2 \, dx$\\被積分関数 $f(x) = 3x^2$ の原始関数の一つは $F(x) = x^3$ である（実際、$F'(x) = 3x^2$）。\\微積分学の基本定理 $\int_a^b f(x)dx = [F(x)]_a^b = F(b) - F(a)$ を適用すると、\\$S = [x^3]_1^3 = 3^3 - 1^3 = 27 - 1 = 26$\\アルキメデスが古代に膨大な取り尽くし法で求めた放物線下の面積が、微積分学の基本定理によってわずか2行で厳密に計算できる。
\end{examplebox}


\section*{第12章の確認演習}
\addcontentsline{toc}{section}{第12章の確認演習}
\begin{enumerate}[label=\textbf{問\arabic*.}]
\item 区分求積法（リーマン和）において、短冊の幅をゼロに近づけたとき定積分へと収束するメカニズムを解説せよ。
\item 微積分学の基本定理 $\frac{d}{dx}\left(\int_a^x f(t)dt\right) = f(x)$ を幾何学的に証明せよ。
\item 現代の人工知能（ニューラルネットワークの勾配降下法）において導関数がどのように活用されているか論ぜよ。
\end{enumerate}

\chapter*{付録：高校生のための微積分重要用語辞典}
\addcontentsline{toc}{chapter}{付録：高校生のための微積分重要用語辞典}
\markboth{付録：微積分重要用語辞典}{付録：微積分重要用語辞典}

本書で登場した極めて重要かつ本質的な数学・物理の用語を、トンプソンの直観的視点から整理した解説集である。

\begin{description}[leftmargin=3.5cm, style=nextline]
\item[微小量（Differential, $dx, dy$）]
ゼロではないが、あらゆる通常の測定可能な有限量と比べて極めて微小な増分。トンプソンはこれを「ほんの少しのかけら（a little bit of）」と表現した。

\item[小ささの階級（Degrees of Smallness）]
1次微小量 $dx$ に対し、2次微小量 $(dx)^2$ や $dx \cdot dy$ は桁違いに小さく、極限操作において正当に無視できるという近代微積分の近似原理。

\item[導関数（Derivative, $\frac{dy}{dx}$）]
変数 $x$ の微小な成長に対する変数 $y$ の微小な成長の比率（相対的成長度）。グラフ上ではその点における接線の傾きを表す。

\item[平均変化率（Average Rate of Change）]
区間 $[x, x + \Delta x]$ における関数の平均的な変化の傾き $\frac{\Delta y}{\Delta x} = \frac{f(x+\Delta x) - f(x)}{\Delta x}$。

\item[瞬間変化率（Instantaneous Rate of Change）]
平均変化率の区間幅 $\Delta x$ を限りなくゼロに近づけた極限値。スピードメーターが指すその瞬間の速度。

\item[接線（Tangent Line）]
曲線上の点において、曲線の局所的な進行方向と完全に一致する直線。傾きは導関数 $f'(x)$ に等しい。

\item[ベキ関数の微分公式]
自然数、負数、分数を含む任意の実数指数 $n$ に対し、$\frac{d}{dx}(x^n) = n x^{n-1}$ が成立するという微積分の基本定理。

\item[積の微分公式（Product Rule）]
2つの関数 $u, v$ の積の導関数は $(uv)' = u'v + uv'$ となる。長方形の面積拡大モデルにより直観的に証明される。

\item[商の微分公式（Quotient Rule）]
分数の導関数は $\left(\frac{u}{v}\right)' = \frac{u'v - uv'}{v^2}$ となる。分母の増大が商を減少させるため分子にマイナス符号が現れる。

\item[連鎖律（Chain Rule）]
合成関数 $y = f(g(x))$ に対し、中間変数 $u = g(x)$ を介して $\frac{dy}{dx} = \frac{dy}{du} \cdot \frac{du}{dx}$ と分解・計算する強力な技法（トンプソンの「便利な裏技」）。

\item[第2次導関数（Second Derivative, $\frac{d^2y}{dx^2}$）]
導関数をさらに微分したもの。曲線の曲がり具合（凹凸）や、物理における加速度（速度の変化率）を表す。

\item[極大値・極小値（Maxima and Minima）]
関数が局所的に最大または最小となる値。その地点では接線が水平（$\frac{dy}{dx} = 0$）となり、2次導関数の符号によって山（$< 0$）か谷（$> 0$）かが判定される。

\item[変曲点（Inflection Point）]
曲線の凹凸（上に凸と下に凸）が入れ替わる境界点。第2次導関数がゼロ（$\frac{d^2y}{dx^2} = 0$）となる。

\item[ネイピア数（Napier's Constant, $e$）]
連続複利の極限 $\lim_{n \to \infty} \left(1 + \frac{1}{n}\right)^n = 2.71828\dots$ で定義される超越数。関数 $y = e^x$ は微分しても自分自身のままである唯一無二の性質を持つ。

\item[自然対数（Natural Logarithm, $\ln x$）]
ネイピア数 $e$ を底とする対数。その導関数は $\frac{d}{dx}(\ln x) = \frac{1}{x}$ となり、反比例関数の原始関数を与える。

\item[弧度法（Radian）]
半径と弧長の比によって角度を定義する単位系。三角関数の極限 $\lim_{\theta \to 0} \frac{\sin\theta}{\theta} = 1$ を保証し、微積分を美しく成立させる。

\item[不定積分（Indefinite Integral, $\int f(x) dx$）]
微分の逆演算として、導関数 $f(x)$ から元の関数（原始関数）$F(x) + C$ を復元する操作。

\item[積分定数（Constant of Integration, $C$）]
定数は微分するとゼロになるため、積分によって復元する際に生じる任意定数。初期条件によって特定される。

\item[定積分（Definite Integral, $\int_a^b f(x) dx$）]
区間 $[a, b]$ における微小短冊 $f(x)dx$ の無限和。曲線と $x$ 軸で囲まれた符号付き面積を表す。

\item[微積分学の基本定理（Fundamental Theorem of Calculus）]
微分と積分が表裏一体の逆演算であり、面積関数 $S(x)$ の導関数が元の関数の高さ $f(x)$ に等しいという大定理：$\int_a^b f(x)dx = F(b) - F(a)$。
\end{description}

\chapter*{付録：高校数学総合演習問題と詳細解説}
\addcontentsline{toc}{chapter}{付録：高校数学総合演習問題と詳細解説}
\markboth{付録：高校数学総合演習問題}{付録：高校数学総合演習問題}

全12章で学んだトンプソン流微積分の理論を総合的に活用する、大学入試・実践レベルの演習問題である。

\section*{【総合演習問題】}

\begin{enumerate}[label=\textbf{第\arabic*問.}]
\item 関数 $f(x) = x^3 - 6x^2 + 9x + 2$ について、次の問いに答えよ。
\begin{enumerate}
\item 導関数 $f'(x)$ および第2次導関数 $f''(x)$ を求めよ。
\item 増減表を作成し、極大値・極小値および変曲点の座標を求めよ。
\item $x \ge 0$ における曲線 $y = f(x)$ の概形を描け。
\end{enumerate}

\item 半径 $R$ の球に内接する直円柱を考える。
\begin{enumerate}
\item 円柱の底面の半径を $r$、高さを $2h$ とするとき、$r, h, R$ の関係式を求めよ。
\item 円柱の体積 $V$ を $h$ の関数として表せ。
\item 円柱の体積 $V$ が最大となるときの高さ $2h$ および最大体積 $V_{\text{max}}$ を $R$ を用いて表せ。
\end{enumerate}

\item 次の不定積分および定積分を計算せよ。
\begin{enumerate}
\item $\int (3x^2 - 4x + 1)\,dx$
\item $\int \left(e^{2x} + \frac{1}{x}\right)dx \quad (x > 0)$
\item $\int_0^{\pi} \sin(x)\,dx$
\item $\int_1^2 \left(x + \frac{1}{x^2}\right)dx$
\end{enumerate}

\item 放物線 $y = 4 - x^2$ と直線 $y = x + 2$ について、次の問いに答えよ。
\begin{enumerate}
\item 放物線と直線の交点の $x$ 座標を求めよ。
\item 放物線と直線で囲まれた図形の面積 $S$ を定積分を用いて求めよ。
\end{enumerate}

\item 地上から初速度 $v_0$、仰角 $\theta$ で打ち出された物体の運動を考える（空気抵抗は無視し、重力加速度を $g$ とする）。
\begin{enumerate}
\item 水平方向の位置 $x(t)$ および鉛直方向の位置 $y(t)$ を時間の関数として表せ。
\item 軌道方程式 $y = f(x)$ を求めよ。
\item 水平到達距離 $X$ が最大となる仰角 $\theta$ を、微分の考え方を用いて導け。
\end{enumerate}
\end{enumerate}

\newpage
\section*{【総合演習問題の解答・解説】}

\subsection*{第1問の解答}
(a) $f(x) = x^3 - 6x^2 + 9x + 2$ に対し、\\
$f'(x) = 3x^2 - 12x + 9 = 3(x^2 - 4x + 3) = 3(x - 1)(x - 3)$。\\
$f''(x) = 6x - 12 = 6(x - 2)$。\\
(b) $f'(x) = 0$ となるのは $x = 1, 3$。\\
$x = 1$ のとき、$f(1) = 1 - 6 + 9 + 2 = 6$。$f''(1) = -6 < 0$ より極大値 $6$。\\
$x = 3$ のとき、$f(3) = 27 - 54 + 27 + 2 = 2$。$f''(3) = 6 > 0$ より極小値 $2$。\\
変曲点は $f''(x) = 0$ となる $x = 2$ であり、座標は $(2, f(2)) = (2, 4)$ である。\\
(c) 曲線は $(0, 2)$ から出発し、$(1, 6)$ で極大となり、$(2, 4)$ で凹凸が反転し、$(3, 2)$ で極小となった後、右肩上がりに急上昇する三次曲線となる。

\subsection*{第2問の解答}
(a) 球の中心を原点にとると、円柱の上底面の縁は球面上にあるため、三平方の定理より $r^2 + h^2 = R^2$、すなわち $r^2 = R^2 - h^2$（ただし $0 < h < R$）。\\
(b) 円柱の底面積は $\pi r^2 = \pi(R^2 - h^2)$、高さは $2h$ であるから、体積は\\
$V(h) = \pi(R^2 - h^2)(2h) = 2\pi(R^2 h - h^3)$。\\
(c) $V(h)$ を $h$ で微分して極大値を求める：\\
$\frac{dV}{dh} = 2\pi(R^2 - 3h^2)$。\\
$\frac{dV}{dh} = 0$ とおくと、$h > 0$ より $h = \frac{R}{\sqrt{3}} = \frac{\sqrt{3}}{3}R$。\\
$\frac{d^2V}{dh^2} = -12\pi h < 0$ より、このとき体積は最大となる。\\
求める高さは $2h = \frac{2\sqrt{3}}{3}R$。\\
最大体積は $V_{\text{max}} = 2\pi\left(R^2 \cdot \frac{R}{\sqrt{3}} - \frac{R^3}{3\sqrt{3}}\right) = 2\pi \cdot \frac{2R^3}{3\sqrt{3}} = \frac{4\sqrt{3}\pi}{9}R^3$ である。

\subsection*{第3問の解答}
(a) $\int (3x^2 - 4x + 1)dx = x^3 - 2x^2 + x + C$。\\
(b) 連鎖律の逆より $\int e^{2x}dx = \frac{1}{2}e^{2x}$。したがって $\frac{1}{2}e^{2x} + \ln|x| + C$。\\
(c) $[-\cos(x)]_0^{\pi} = (-\cos\pi) - (-\cos 0) = -(-1) - (-1) = 1 + 1 = 2$。\\
(d) $\left[\frac{x^2}{2} - \frac{1}{x}\right]_1^2 = \left(2 - \frac{1}{2}\right) - \left(\frac{1}{2} - 1\right) = \frac{3}{2} - \left(-\frac{1}{2}\right) = 2$。

\subsection*{第4問の解答}
(a) 交点の方程式は $4 - x^2 = x + 2 \iff x^2 + x - 2 = 0 \iff (x + 2)(x - 1) = 0$。\\
交点の $x$ 座標は $x = -2, 1$。\\
(b) 区間 $[-2, 1]$ において放物線が直線の上側にある（$4 - x^2 \ge x + 2$）。\\
面積 $S = \int_{-2}^1 \{(4 - x^2) - (x + 2)\}dx = \int_{-2}^1 (-x^2 - x + 2)dx$。\\
これは有名な $\frac{1}{6}$ 公式 $S = \frac{1}{6}(\beta - \alpha)^3$ の形である：\\
$S = \frac{1}{6}(1 - (-2))^3 = \frac{1}{6}(3)^3 = \frac{27}{6} = \frac{9}{2} = 4.5$。

\subsection*{第5問の解答}
(a) 水平方向は等速直線運動：$x(t) = (v_0 \cos\theta)t$。\\
鉛直方向は等加速度運動：$y(t) = (v_0 \sin\theta)t - \frac{1}{2}gt^2$。\\
(b) $t = \frac{x}{v_0 \cos\theta}$ を $y(t)$ に代入すると、軌道方程式は\\
$y = (\tan\theta)x - \frac{g}{2v_0^2 \cos^2\theta}x^2$。\\
(c) 再び地面（$y = 0$）に落下する時刻は $t = \frac{2v_0 \sin\theta}{g}$。\\
このときの水平到達距離は\\
$X(\theta) = (v_0 \cos\theta) \cdot \frac{2v_0 \sin\theta}{g} = \frac{v_0^2}{g}(2\sin\theta\cos\theta) = \frac{v_0^2}{g}\sin(2\theta)$。\\
$X(\theta)$ を $\theta$ で微分すると $\frac{dX}{d\theta} = \frac{2v_0^2}{g}\cos(2\theta)$。\\
$\frac{dX}{d\theta} = 0$ とおくと $2\theta = 90^\circ \implies \theta = 45^\circ$。\\
$\theta = 45^\circ$ のとき水平到達距離は最大値 $X_{\text{max}} = \frac{v_0^2}{g}$ となる。

\end{document}
